\documentclass[11pt,a4paper]{article}

\usepackage[T1]{fontenc}
\usepackage[utf8]{inputenc}
\usepackage{lmodern}
\usepackage[margin=2.2cm,top=2.4cm,bottom=2.4cm]{geometry}
\usepackage{xcolor}
\usepackage{listings}
\usepackage{enumitem}
\usepackage{fancyhdr}
\usepackage{titlesec}
\usepackage{amsmath}
\usepackage{amssymb}
\usepackage{booktabs}
\usepackage{array}
\usepackage{tikz}
\usetikzlibrary{arrows.meta,positioning,calc,fit,backgrounds,decorations.pathreplacing}
\usepackage[hidelinks,bookmarks,bookmarksnumbered]{hyperref}

% ----------------------------------------------------------------- colours
\definecolor{ink}{HTML}{15171C}
\definecolor{accent}{HTML}{1D4ED8}
\definecolor{muted}{HTML}{5C6370}
\definecolor{rule}{HTML}{D3D7DE}
\definecolor{codebg}{HTML}{F6F7F9}
\definecolor{kw}{HTML}{0B5FA5}
\definecolor{cm}{HTML}{6A7380}
\definecolor{st}{HTML}{9B2C2C}
\definecolor{bad}{HTML}{B42318}
\definecolor{ok}{HTML}{1A7F37}
\definecolor{easybg}{HTML}{E7F6EC}
\definecolor{medbg}{HTML}{FDF3DC}
\definecolor{hardbg}{HTML}{FBE9E7}

\hypersetup{pdftitle={Java Coding Interview - Revision Book},
            pdfauthor={Assessment},
            pdfsubject={60 worked problems for a CodeSignal-style Java screen}}

% ----------------------------------------------------------------- listings
\lstdefinestyle{java}{
  language=Java,
  basicstyle=\ttfamily\footnotesize,
  keywordstyle=\color{kw}\bfseries,
  commentstyle=\color{cm}\itshape,
  stringstyle=\color{st},
  numbers=left,
  numberstyle=\tiny\color{muted},
  numbersep=9pt,
  backgroundcolor=\color{codebg},
  frame=leftline,
  framesep=7pt,
  framerule=1.1pt,
  rulecolor=\color{accent},
  xleftmargin=18pt,
  breaklines=true,
  breakatwhitespace=true,
  showstringspaces=false,
  columns=fullflexible,
  keepspaces=true,
  tabsize=4,
  aboveskip=8pt,
  belowskip=8pt,
}
\lstdefinestyle{plain}{
  basicstyle=\ttfamily\footnotesize,
  backgroundcolor=\color{codebg},
  frame=leftline,
  framesep=7pt,
  framerule=1.1pt,
  rulecolor=\color{rule},
  xleftmargin=10pt,
  breaklines=true,
  showstringspaces=false,
  columns=fullflexible,
  keepspaces=true,
  aboveskip=7pt,
  belowskip=7pt,
}
\lstset{style=java}

% ----------------------------------------------------------------- headings
\titleformat{\section}{\Large\bfseries\color{accent}}{\thesection.}{0.6em}{}
\titlespacing*{\section}{0pt}{0pt}{14pt}
\newcommand{\chapterpage}[2]{\clearpage\section{#1}\noindent\textit{\color{muted}#2}\par\medskip}

\pagestyle{fancy}
\fancyhf{}
\renewcommand{\headrulewidth}{0.4pt}
\fancyhead[L]{\small\color{muted}Java Coding Interview --- Revision Book}
\fancyhead[R]{\small\color{muted}\thepage}

% ----------------------------------------------------------------- problems
\newcounter{probnum}
% NB: \tag is taken by amsmath, hence \dchip.
\newcommand{\dchip}[2]{\colorbox{#1}{\small\textsf{\,#2\,}}}
\newcommand{\difficulty}[1]{%
  \ifnum\pdfstrcmp{#1}{Easy}=0 \dchip{easybg}{Easy}\fi
  \ifnum\pdfstrcmp{#1}{Medium}=0 \dchip{medbg}{Medium}\fi
  \ifnum\pdfstrcmp{#1}{Hard}=0 \dchip{hardbg}{Hard}\fi
}
\newcommand{\problem}[2]{%
  \refstepcounter{probnum}%
  \par\medskip\noindent
  \rule{\linewidth}{0.4pt}\par\smallskip
  {\large\bfseries P\theprobnum.\ #1}\hfill\difficulty{#2}\par\smallskip
  \addcontentsline{toc}{subsection}{P\theprobnum.\ #1}%
  \nopagebreak
}
\newcommand{\solheading}{\par\smallskip\noindent\textbf{\color{accent}Solution}\par\nopagebreak}
\newcommand{\idea}[1]{\par\noindent\textbf{Idea.}\ #1\par}
\newcommand{\cost}[1]{\par\noindent\textbf{Cost.}\ #1\par}
\newcommand{\trap}[1]{\par\smallskip\noindent\textbf{\color{bad}Trap.}\ #1\par}
\newcommand{\alsoknown}[1]{\par\noindent{\small\color{muted}Also appears as: #1}\par}

% ----------------------------------------------------------------- diagrams
\definecolor{dline}{HTML}{BFC6D1}
\definecolor{dplain}{HTML}{FFFFFF}
\definecolor{dsoft}{HTML}{EDF2FA}
\definecolor{dhot}{HTML}{FAE0D4}
\definecolor{dwarm}{HTML}{FBEED5}
\definecolor{dcool}{HTML}{DFF2EA}
\definecolor{ddead}{HTML}{EFEFEA}

\tikzset{
  cell/.style      ={draw=dline, line width=.5pt, fill=dplain, minimum size=7.4mm,
                     inner sep=0pt, font=\ttfamily\small, text=ink},
  soft/.style      ={cell, fill=dsoft},
  hot/.style       ={cell, fill=dhot},
  warm/.style      ={cell, fill=dwarm},
  cool/.style      ={cell, fill=dcool},
  dead/.style      ={cell, fill=ddead, text=muted},
  disc/.style      ={draw=dline, line width=.5pt, fill=dsoft, circle,
                     minimum size=7.4mm, inner sep=0pt, font=\ttfamily\small, text=ink},
  idx/.style       ={font=\scriptsize, text=muted, inner sep=1.5pt},
  ptr/.style       ={font=\footnotesize\bfseries, text=accent, inner sep=1.5pt},
  note/.style      ={font=\footnotesize, text=muted, inner sep=2pt},
  flow/.style      ={-{Stealth[length=4.5pt,width=3.5pt]}, draw=accent, line width=.7pt},
  back/.style      ={-{Stealth[length=4.5pt,width=3.5pt]}, draw=muted, line width=.6pt,
                     densely dashed},
  edge/.style      ={draw=dline, line width=.6pt},
  span/.style      ={decorate, decoration={brace, amplitude=4pt, raise=2pt},
                     draw=accent, line width=.6pt},
}

% A centred diagram block.
\newenvironment{dia}
  {\par\addvspace{9pt}\begin{center}\begin{tikzpicture}}
  {\end{tikzpicture}\end{center}\addvspace{3pt}}

% \cellrow{prefix}{y}{a,b,c}  -- a horizontal strip of cells, nodes <prefix>-0, -1, ...
\newcommand{\cellrow}[3]{%
  \foreach \cv [count=\ci from 0] in {#3} {%
    \node[cell] (#1-\ci) at (\ci*0.8, #2) {\cv};}}

% \cellrowst{prefix}{y}{style/value, style/value, ...}  -- same, per-cell styles
\newcommand{\cellrowst}[3]{%
  \foreach \cs/\cv [count=\ci from 0] in {#3} {%
    \node[\cs] (#1-\ci) at (\ci*0.8, #2) {\cv};}}

% \idxrow{y}{last index}  -- index numbers under a strip
\newcommand{\idxrow}[2]{%
  \foreach \ci in {0,...,#2} {\node[idx] at (\ci*0.8, #1) {\ci};}}

% \caplabel{x}{y}{text}  -- a small grey annotation
\newcommand{\caplabel}[3]{\node[note] at (#1, #2) {#3};}

\emergencystretch=2.5em
\hbadness=2000

\setlist[itemize]{leftmargin=1.2em,itemsep=2pt,topsep=4pt}
\setlist[enumerate]{leftmargin=1.5em,itemsep=2pt,topsep=4pt}

\begin{document}

% ----------------------------------------------------------------- title
\begin{titlepage}
\centering
\vspace*{3.5cm}
{\Huge\bfseries Java Coding Interview\par}
\vspace{0.5cm}
{\Huge\bfseries Revision Book\par}
\vspace{1.2cm}
{\large\color{muted} 60 worked problems for a CodeSignal-style screen\par}
\vspace{0.3cm}
{\large\color{muted} Final-year engineering level\par}
\vspace{2.5cm}
\rule{8cm}{1pt}\par
\vspace{2.5cm}
{\color{muted}
Every problem has a complete Java solution, its complexity,\\[2pt]
and the mistake that most often costs the points.\par}
\vfill
{\small\color{muted} Companion to \texttt{ASSESSMENT.md} --- work the book first, sit the assessment second.\par}
\end{titlepage}

% ----------------------------------------------------------------- front
\setcounter{page}{1}
\section*{How to use this book}
\addcontentsline{toc}{section}{How to use this book}

The screen you are preparing for is reported to be four problems in seventy
minutes on CodeSignal, with tasks in the range of \emph{array manipulation} and
\emph{data processing: calculating averages, filtering data}. That is the bar.
It is not a competitive-programming bar, and almost nothing in rounds like this
needs dynamic programming or graph theory.

What it does need is that you write correct Java \emph{quickly}, without
stopping to remember how a \texttt{HashMap} iterates or what
\texttt{split(",")} does to trailing empty fields. Speed on easy and medium
problems is worth far more than the ability to crack one hard one.

\paragraph{A working method.}
\begin{enumerate}
  \item Read the whole statement before typing. Note every edge case it names:
        \texttt{null}, empty, ties, bounds. They are in the statement because
        they are in the tests.
  \item Say the approach out loud in one sentence. If you cannot, you do not
        have one yet.
  \item Write the happy path. Then go back and add the edge cases you noted.
  \item Re-read the statement against your code before moving on. The most
        expensive bug in a timed test is the one you never look for.
\end{enumerate}

\paragraph{How to drill.} Give yourself the time budget printed on each
difficulty: Easy 8 minutes, Medium 15, Hard 25. Write the solution without
looking. Then read the \textbf{Trap} note --- it is the thing that actually
breaks submissions, and it is worth more than the solution itself.

\paragraph{Chapters 1 to 8 are the exam.} Chapter 9 is deliberately above the
bar. Work it only when the rest is comfortable; it is there so that a hard
problem on the day does not surprise you.

\clearpage
\tableofcontents
\clearpage

\IfFileExists{ch0-toolkit.tex}{\chapterpage{The Java toolkit}{The twenty things you must not have to think
about on the day. Read this once; come back to it whenever a drill makes you
hesitate.}

\subsection*{Choosing the container}
\addcontentsline{toc}{subsection}{Choosing the container}

\begin{tabular}{@{}p{4.3cm}p{4.1cm}p{6.6cm}@{}}
\toprule
\textbf{You need} & \textbf{Use} & \textbf{Why / cost} \\
\midrule
Count things by key & \texttt{HashMap<K,Integer>} & \texttt{merge(k,1,Integer::sum)}, O(1) \\
Group things by key & \texttt{HashMap<K,List<V>>} & \texttt{computeIfAbsent(k, x -> new ArrayList<>())} \\
Seen-before / dedupe & \texttt{HashSet<T>} & \texttt{add} returns \texttt{false} if present \\
Keys in sorted order & \texttt{TreeMap<K,V>} & O(log n), gives \texttt{firstKey}, \texttt{floorKey} \\
Insertion order kept & \texttt{LinkedHashMap} & iteration order = insertion order \\
Growable array & \texttt{ArrayList<T>} & O(1) amortised add, O(1) get \\
Stack or queue & \texttt{ArrayDeque<T>} & \texttt{push/pop}, \texttt{offer/poll}; never \texttt{Stack} \\
Top-k / smallest-first & \texttt{PriorityQueue<T>} & O(log n) per op, min-heap by default \\
\bottomrule
\end{tabular}

\medskip
\noindent\texttt{ArrayDeque} is the right answer for both stacks and queues.
\texttt{java.util.Stack} is a synchronised relic and \texttt{LinkedList} costs
you a pointer chase per element.

\subsection*{Idioms worth memorising}
\addcontentsline{toc}{subsection}{Idioms worth memorising}

\begin{lstlisting}[style=java,numbers=none]
// count occurrences
Map<Character,Integer> freq = new HashMap<>();
for (char c : s.toCharArray()) freq.merge(c, 1, Integer::sum);

// group into lists
Map<String,List<String>> byKey = new HashMap<>();
byKey.computeIfAbsent(key, k -> new ArrayList<>()).add(value);

// accumulate a total and a count together
map.computeIfAbsent(dept, k -> new long[2])[0] += salary;   // [0]=total
map.get(dept)[1]++;                                          // [1]=count

// sort with a tie-break: primary descending, secondary ascending
list.sort((a, b) -> a.score != b.score
        ? Integer.compare(b.score, a.score)
        : a.name.compareTo(b.name));

// iterate a map
for (Map.Entry<String,Integer> e : map.entrySet())
    use(e.getKey(), e.getValue());

// array <-> list
int[] arr = list.stream().mapToInt(Integer::intValue).toArray();
String[] out = list.toArray(new String[0]);

// 2D array, all cells set
char[][] g = new char[rows][cols];
for (char[] row : g) Arrays.fill(row, '.');

// the four neighbours (eight: nest -1..1 and skip 0,0)
int[][] DIRS = {{-1,0},{1,0},{0,-1},{0,1}};
\end{lstlisting}

\subsection*{Traps that cost real points}
\addcontentsline{toc}{subsection}{Traps that cost real points}

\begin{itemize}
  \item \textbf{\texttt{split(",")} drops trailing empties.} \texttt{"a,b,"}
        splits into 2 parts, not 3. Use \texttt{split(",", -1)} whenever the
        field count is part of your validation.
  \item \textbf{\texttt{\%} keeps the sign of the dividend.}
        \texttt{-3 \% 2 == -1}, so \texttt{n \% 2 == 1} is false for negative
        odd numbers. Write \texttt{n \% 2 != 0}.
  \item \textbf{Integer division truncates toward zero.} \texttt{7/2 == 3} and
        \texttt{-7/2 == -3}. That is not floor for negatives.
  \item \textbf{\texttt{int} overflow is silent.} Sums of up to
        $2\times10^9$ values overflow. Accumulate in \texttt{long}.
  \item \textbf{\texttt{==} on boxed types compares references.} Use
        \texttt{.equals}, or unbox to \texttt{int} first. \texttt{Integer}
        caches only $-128$ to $127$, which is why small cases seem to work.
  \item \textbf{\texttt{String +=} in a loop is O(n$^2$).} Use
        \texttt{StringBuilder}.
  \item \textbf{\texttt{Concurrent\-Modification\-Exception}} when you add to or
        remove from a collection while a for-each is running. Collect the keys
        first, then mutate; or use \texttt{Iterator.remove()}.
  \item \textbf{A recursive flood fill overflows the stack} at a few tens of
        thousands of cells. Use an explicit \texttt{ArrayDeque}.
  \item \textbf{\texttt{Arrays.sort} on primitives is a dual-pivot quicksort}
        with an O(n$^2$) worst case on crafted input; on objects it is a stable
        merge sort. Boxing to \texttt{Integer[]} is the paranoid fix, and
        almost never needed.
  \item \textbf{Returning \texttt{null} where an empty array was asked for.}
        Read the return contract twice.
\end{itemize}

\subsection*{Complexity you should recognise instantly}
\addcontentsline{toc}{subsection}{Complexity you should recognise instantly}

\begin{tabular}{@{}llp{6.4cm}@{}}
\toprule
\textbf{Shape} & \textbf{Cost} & \textbf{Safe up to} \\
\midrule
Single pass & $O(n)$ & $10^8$ \\
Sort & $O(n\log n)$ & $10^7$ \\
Nested loop over the input & $O(n^2)$ & $\sim\!10^4$ \\
Pass per element, each a scan & $O(n^2)$ & the same, and this is the one you write by accident \\
Flood fill / BFS on a grid & $O(\text{rows}\times\text{cols})$ & $10^7$ cells \\
\bottomrule
\end{tabular}

\medskip
\noindent If the constraints say $n \le 200\,000$, an $O(n^2)$ solution will
time out. That number in a statement is the examiner telling you the intended
complexity.
}{}
\IfFileExists{ch1-strings.tex}{\chapterpage{Strings}{Eight problems. Strings are the most common opener in a
screen because they are quick to state and full of edge cases.}

% ------------------------------------------------------------------
\problem{Reverse the words in a sentence}{Easy}

Given a sentence, return it with the words in reverse order and exactly one
space between them. Leading, trailing and repeated spaces are removed.
\texttt{"\ the sky\ \ is blue "} becomes \texttt{"blue is sky the"}.

\idea{Split on runs of whitespace, walk the parts backwards, join with a single space.}

\solheading
\begin{lstlisting}[style=java]
public static String reverseWords(String s) {
    if (s == null) return "";
    String[] parts = s.trim().split("\\s+");
    StringBuilder out = new StringBuilder();
    for (int i = parts.length - 1; i >= 0; i--) {
        if (parts[i].isEmpty()) continue;
        if (out.length() > 0) out.append(' ');
        out.append(parts[i]);
    }
    return out.toString();
}
\end{lstlisting}

\cost{$O(n)$ time, $O(n)$ space.}
\trap{Splitting on \texttt{" "} instead of \texttt{"\textbackslash\textbackslash s+"} leaves empty strings in the
array when spaces repeat. And \texttt{trim()} first, or a leading space gives
you an empty first element anyway.}

% ------------------------------------------------------------------
\problem{Valid palindrome, ignoring punctuation}{Easy}

Return \texttt{true} if the string reads the same forwards and backwards when
you consider only letters and digits, and ignore case.
\texttt{"A man, a plan, a canal: Panama"} is a palindrome.

\idea{Two pointers walking inwards, each skipping anything that is not alphanumeric.}

\solheading
\begin{lstlisting}[style=java]
public static boolean isPalindrome(String s) {
    if (s == null) return true;
    int i = 0, j = s.length() - 1;
    while (i < j) {
        while (i < j && !Character.isLetterOrDigit(s.charAt(i))) i++;
        while (i < j && !Character.isLetterOrDigit(s.charAt(j))) j--;
        if (Character.toLowerCase(s.charAt(i)) !=
            Character.toLowerCase(s.charAt(j))) return false;
        i++; j--;
    }
    return true;
}
\end{lstlisting}

\cost{$O(n)$ time, $O(1)$ space.}
\trap{The inner skip loops need their own \texttt{i < j} guard, or a string of
pure punctuation runs the pointer off the end.}

% ------------------------------------------------------------------
\problem{First non-repeating character}{Easy}

Return the index of the first character that appears exactly once, or
$-1$. In \texttt{"loveleetcode"} the answer is 2.

\idea{Count every character in one pass, then scan the string again in order and
return the first with count 1. Two passes beat one clever pass.}

\solheading
\begin{lstlisting}[style=java]
public static int firstUnique(String s) {
    if (s == null || s.isEmpty()) return -1;
    int[] freq = new int[128];                 // ASCII
    for (int i = 0; i < s.length(); i++) freq[s.charAt(i)]++;
    for (int i = 0; i < s.length(); i++)
        if (freq[s.charAt(i)] == 1) return i;
    return -1;
}
\end{lstlisting}

\cost{$O(n)$ time, $O(1)$ space (the table is a fixed 128).}
\trap{Scanning the \emph{map} rather than the string for the second pass loses
the ordering. \texttt{HashMap} has no order; if you must use one, use
\texttt{LinkedHashMap}, or just re-scan the string as above.}

% ------------------------------------------------------------------
\problem{Run-length encoding}{Easy}

Replace each maximal run of identical characters with the character followed by
the run length, writing a run of one as just the character.
\texttt{"aaabccddd"} becomes \texttt{"a3bc2d3"}.

\idea{One pointer at the start of the run, a second scanning forward while the
character is unchanged.}

\solheading
\begin{lstlisting}[style=java]
public static String compress(String s) {
    if (s == null || s.isEmpty()) return "";
    StringBuilder out = new StringBuilder();
    int i = 0;
    while (i < s.length()) {
        char ch = s.charAt(i);
        int j = i;
        while (j < s.length() && s.charAt(j) == ch) j++;
        out.append(ch);
        if (j - i > 1) out.append(j - i);
        i = j;
    }
    return out.toString();
}
\end{lstlisting}

\cost{$O(n)$ time, $O(n)$ space.}
\trap{Run lengths can have more than one digit --- \texttt{"z12"} is a valid
output. Any solution that appends a single \texttt{char} for the count is
wrong past nine.}
\alsoknown{Task 1 of the assessment.}

% ------------------------------------------------------------------
\problem{Anagram check}{Easy}

Return \texttt{true} if two strings contain the same letters with the same
multiplicities.

\idea{Count up for the first string, count down for the second, and check nothing
is left over. Length differing is an instant \texttt{false}.}

\solheading
\begin{lstlisting}[style=java]
public static boolean isAnagram(String a, String b) {
    if (a == null || b == null || a.length() != b.length()) return false;
    int[] freq = new int[26];
    for (int i = 0; i < a.length(); i++) {
        freq[a.charAt(i) - 'a']++;
        freq[b.charAt(i) - 'a']--;
    }
    for (int f : freq) if (f != 0) return false;
    return true;
}
\end{lstlisting}

\cost{$O(n)$ time, $O(1)$ space.}
\trap{The \texttt{- 'a'} indexing assumes lowercase letters only. If the input
can be any character, use a \texttt{HashMap} or a 128-entry table --- an
uppercase letter would write to a negative index and throw.}

% ------------------------------------------------------------------
\problem{Longest common prefix}{Easy}

Given an array of strings, return the longest prefix common to all of them, or
\texttt{""}.

\idea{Compare column by column: take character $i$ of the first string and check
it against character $i$ of every other. Stop at the first mismatch or the
first string that ends.}

\solheading
\begin{lstlisting}[style=java]
public static String longestCommonPrefix(String[] words) {
    if (words == null || words.length == 0) return "";
    for (int i = 0; i < words[0].length(); i++) {
        char c = words[0].charAt(i);
        for (int w = 1; w < words.length; w++) {
            if (words[w] == null) return "";
            if (i >= words[w].length() || words[w].charAt(i) != c)
                return words[0].substring(0, i);
        }
    }
    return words[0];
}
\end{lstlisting}

\cost{$O(\text{total characters})$ time, $O(1)$ extra space.}
\trap{Forgetting the \texttt{i >= words[w].length()} check throws
\texttt{StringIndexOutOfBounds} as soon as one word is a strict prefix of
another (\texttt{["ab", "a"]}).}

% ------------------------------------------------------------------
\problem{Compare two version strings}{Medium}

Versions are dot-separated numbers: \texttt{"1.2.10"}. Return $-1$, $0$ or $1$
for less than, equal, greater than. Missing components count as zero, so
\texttt{"1.0"} equals \texttt{"1"}.

\idea{Split both, then walk to the longer length, treating a missing component as
zero. Compare numerically, never as text.}

\solheading
\begin{lstlisting}[style=java]
public static int compareVersions(String a, String b) {
    String[] x = a.split("\\.");
    String[] y = b.split("\\.");
    int n = Math.max(x.length, y.length);
    for (int i = 0; i < n; i++) {
        long left  = i < x.length ? Long.parseLong(x[i]) : 0;
        long right = i < y.length ? Long.parseLong(y[i]) : 0;
        if (left != right) return left < right ? -1 : 1;
    }
    return 0;
}
\end{lstlisting}

\cost{$O(n)$ time.}
\trap{Two of them. \texttt{split(".")} splits on \emph{every character},
because \texttt{.} is a regex wildcard --- you must escape it. And comparing
the components as strings makes \texttt{"1.10"} less than \texttt{"1.9"}.}

% ------------------------------------------------------------------
\problem{Longest substring without repeating characters}{Medium}

Return the length of the longest substring in which no character repeats.
In \texttt{"abcabcbb"} the answer is 3.

\idea{A sliding window. Remember the last index at which each character was seen;
when you meet a repeat that is inside the window, jump the left edge past it.}

\solheading
\begin{lstlisting}[style=java]
public static int longestUnique(String s) {
    if (s == null || s.isEmpty()) return 0;
    Map<Character,Integer> lastSeen = new HashMap<>();
    int best = 0, left = 0;
    for (int right = 0; right < s.length(); right++) {
        char c = s.charAt(right);
        Integer prev = lastSeen.get(c);
        if (prev != null && prev >= left) left = prev + 1;
        lastSeen.put(c, right);
        best = Math.max(best, right - left + 1);
    }
    return best;
}
\end{lstlisting}

\cost{$O(n)$ time, $O(\min(n,\Sigma))$ space.}
\trap{The \texttt{prev >= left} test is the whole problem. Without it a
character seen long ago, already outside the window, drags the left edge
\emph{backwards} and the window stops being valid.}
}{}
\IfFileExists{ch2-arrays.tex}{\chapterpage{Arrays and two pointers}{Eight problems. A read index and a write
index, or two ends closing in --- this one idea turns a great many quadratic
solutions linear.}

% ------------------------------------------------------------------
\problem{Move zeros to the end}{Easy}

Move every zero to the end, keeping the order of the non-zero elements, in
place.

\idea{A write pointer. Copy every non-zero forward to the write position, then
fill the tail with zeros.}

\begin{dia}
  \cellrowst{a}{0}{dead/0,soft/1,dead/0,soft/3,soft/12}
  \idxrow{-0.62}{4}
  \node[ptr] at (0.8,0.78) {write};
  \node[ptr] at (2.4,0.78) {read};
  \draw[flow] (0.8,0.62) -- (0.8,0.44);
  \draw[flow] (2.4,0.62) -- (2.4,0.44);
  \draw[flow] (2.4,1.15) to[out=160,in=20] (0.9,1.15);
  \caplabel{5.0}{0.78}{read runs ahead; write only advances on a non-zero};
  \cellrowst{b}{-1.7}{hot/1,hot/3,hot/12,dead/0,dead/0}
  \node[note, anchor=east] at (-0.5,-1.7) {after};
\end{dia}

\textbf{How it works.} Two indices move at different speeds over the same
array. \texttt{read} visits every slot; \texttt{write} only advances when
something is actually kept. Because \texttt{write} can never get ahead of
\texttt{read}, copying \texttt{a[read]} into \texttt{a[write]} can only ever
overwrite a slot that has already been read --- which is why this is safe to
do in place with no temporary array.

When the first loop ends, \texttt{write} is the count of non-zero elements, so
everything from there to the end is tail that must be zeroed.

\solheading
\begin{lstlisting}[style=java]
public static void moveZeros(int[] a) {
    if (a == null) return;
    int write = 0;
    for (int read = 0; read < a.length; read++)
        if (a[read] != 0) a[write++] = a[read];
    while (write < a.length) a[write++] = 0;
}
\end{lstlisting}

\cost{$O(n)$ time, $O(1)$ space.}
\trap{Swapping with the \emph{end} of the array instead of copying to
\texttt{write} destroys the relative order of the non-zero elements, which the
statement asks you to preserve.}

% ------------------------------------------------------------------
\problem{Remove duplicates from a sorted array}{Easy}

Compact a sorted array in place so each value appears once; return the new
length.

\idea{The same write pointer. Because the array is sorted, duplicates are
always adjacent, so compare each element with the last one you kept.}

\begin{dia}
  \cellrowst{a}{0}{soft/1,dead/1,soft/2,dead/2,dead/2,soft/3}
  \idxrow{-0.62}{5}
  \node[ptr] at (1.6,0.78) {write};
  \node[ptr] at (3.2,0.78) {read};
  \draw[flow] (1.6,0.62) -- (1.6,0.44);
  \draw[flow] (3.2,0.62) -- (3.2,0.44);
  \cellrowst{b}{-1.7}{hot/1,hot/2,hot/3,dead/2,dead/2,dead/3}
  \draw[span] (-0.4,-2.3) -- (1.2,-2.3);
  \caplabel{0.4}{-2.85}{returned length 3; the rest is ignored};
\end{dia}

\textbf{How it works. }Sortedness is doing the heavy lifting: it guarantees
that if a value is going to repeat, it repeats immediately. So the test for
``have I seen this already'' collapses from a set lookup to a single
comparison with the previous kept element.

That comparison must be against \texttt{a[write - 1]}, not
\texttt{a[read - 1]}. Once you start overwriting, \texttt{a[read - 1]} may be
a value that was copied there this pass, so it no longer tells you what the
last distinct value was.

\solheading
\begin{lstlisting}[style=java]
public static int dedupeSorted(int[] a) {
    if (a == null || a.length == 0) return 0;
    int write = 1;
    for (int read = 1; read < a.length; read++)
        if (a[read] != a[write - 1]) a[write++] = a[read];
    return write;
}
\end{lstlisting}

\cost{$O(n)$ time, $O(1)$ space.}
\trap{Comparing against \texttt{a[read - 1]}. It looks equivalent and passes
on inputs with no more than two consecutive duplicates.}

% ------------------------------------------------------------------
\problem{Two-sum in a sorted array}{Easy}

Return the indices of two values that add to a target, or \texttt{null}.

\idea{One pointer at each end. If the sum is too big the right pointer must come
in; if it is too small the left must go out. No other move can help.}

\begin{dia}
  \cellrowst{a}{0}{soft/2,cell/7,cell/11,cell/15,soft/19}
  \idxrow{-0.62}{4}
  \node[ptr] at (0,0.78) {i};
  \node[ptr] at (3.2,0.78) {j};
  \draw[flow] (0,0.62) -- (0,0.44);
  \draw[flow] (3.2,0.62) -- (3.2,0.44);
  \node[note, anchor=west] at (3.9,0) {$2+19=21 > 18$};
  \node[note, anchor=west] at (3.9,-1.0) {so \texttt{j--}: nothing paired};
  \node[note, anchor=west] at (3.9,-1.5) {with 19 can be small enough};
\end{dia}

\textbf{How it works.} The argument for why nothing is missed matters more
than the code. Suppose the sum at the current pair is too large. Every pair
using the right-hand value and \emph{any} remaining left-hand value is at
least this large, because the array is sorted ascending and the left pointer
only ever moves right. So that right-hand value cannot be part of any
solution, and discarding it loses nothing. The symmetric argument handles the
too-small case.

Each step eliminates one element, so the search finishes in at most $n$ steps
--- the same answer a hash map would give, with no extra memory.

\solheading
\begin{lstlisting}[style=java]
public static int[] twoSumSorted(int[] a, int target) {
    if (a == null) return null;
    int i = 0, j = a.length - 1;
    while (i < j) {
        long sum = (long) a[i] + a[j];
        if (sum == target) return new int[] {i, j};
        if (sum < target) i++; else j--;
    }
    return null;
}
\end{lstlisting}

On an \emph{unsorted} array this does not apply --- use a \texttt{HashMap}
from value to index instead, also $O(n)$.

\cost{$O(n)$ time, $O(1)$ space.}
\trap{\texttt{a[i] + a[j]} overflows \texttt{int} near the limits. Widen to
\texttt{long} before adding, as above.}

% ------------------------------------------------------------------
\problem{Maximum subarray sum}{Medium}

Return the largest sum obtainable from a contiguous, non-empty subarray.

\noindent\texttt{[-2,1,-3,4,-1,2,1,-5,4]} gives 6.

\idea{Kadane. At each element the best subarray ending here is either this
element alone, or this element extending the best one ending just before it.}

\begin{dia}
  \cellrowst{a}{0}{cell/-2,cell/1,cell/-3,hot/4,hot/-1,hot/2,hot/1,cell/-5,cell/4}
  \idxrow{-0.62}{8}
  \node[note, anchor=east] at (-0.5,0) {input};
  \cellrowst{e}{-1.5}{cell/-2,cell/1,cell/-2,soft/4,soft/3,soft/5,soft/6,cell/1,cell/5}
  \node[note, anchor=east] at (-0.5,-1.5) {best ending here};
  \draw[flow] (4.8,-1.15) -- (4.8,-0.9);
  \caplabel{4.8}{-2.15}{maximum is 6};
\end{dia}

\textbf{How it works.} The second row is the whole algorithm. At each index
you only ever ask one question: is it better to start fresh here, or to carry
forward what came before? If the running total so far is negative it can only
hurt, so you drop it --- that is exactly what
\texttt{Math.max(a[i], endingHere + a[i])} encodes.

Notice that the best overall answer is not necessarily the last value of the
running total; it is the maximum ever reached. That is why \texttt{best} is
tracked separately.

\solheading
\begin{lstlisting}[style=java]
public static int maxSubarraySum(int[] a) {
    if (a == null || a.length == 0) return 0;
    long best = a[0], endingHere = a[0];
    for (int i = 1; i < a.length; i++) {
        endingHere = Math.max(a[i], endingHere + a[i]);
        best = Math.max(best, endingHere);
    }
    return (int) best;
}
\end{lstlisting}

\cost{$O(n)$ time, $O(1)$ space.}
\trap{Initialising \texttt{best} to \texttt{0} returns 0 on an all-negative
array, when the answer is the least negative element. The statement says
non-empty, so the empty subarray is not an option.}

% ------------------------------------------------------------------
\problem{Rotate an array by k}{Medium}

Shift every element right by $k$ positions, wrapping around, in place.

\idea{Reverse the whole array, then reverse the first $k$ and the rest
separately.}

\begin{dia}
  \node[note, anchor=east] at (-0.5,0.9) {start};
  \cellrowst{a}{0.9}{cell/1,cell/2,cell/3,cell/4,cell/5}
  \node[note, anchor=east] at (-0.5,0) {reverse all};
  \cellrowst{b}{0}{soft/5,soft/4,soft/3,soft/2,soft/1}
  \node[note, anchor=east] at (-0.5,-0.9) {reverse each part};
  \cellrowst{c}{-0.9}{hot/4,hot/5,warm/1,warm/2,warm/3}
  \draw[span] (-0.4,-1.5) -- (1.2,-1.5);
  \draw[span] (1.4,-1.5) -- (3.6,-1.5);
  \caplabel{0.4}{-2.0}{first $k=2$};
  \caplabel{2.5}{-2.0}{the rest};
\end{dia}

\textbf{How it works.} Rotating right by $k$ means the last $k$ elements move
to the front, keeping their order, and the rest follow, keeping theirs.
Reversing the whole array brings the last $k$ to the front but in the wrong
order, and does the same to the rest. Reversing each of those two blocks
individually puts both back in order. Three linear reversals, no second
array.

The normalisation \texttt{k = ((k \% n) + n) \% n} handles two cases at once:
$k$ larger than the array, and negative $k$. A single \texttt{\%} leaves a
negative result negative.

\solheading
\begin{lstlisting}[style=java]
public static void rotate(int[] a, int k) {
    if (a == null || a.length == 0) return;
    k = ((k % a.length) + a.length) % a.length;   // handles negative k
    reverse(a, 0, a.length - 1);
    reverse(a, 0, k - 1);
    reverse(a, k, a.length - 1);
}

private static void reverse(int[] a, int i, int j) {
    while (i < j) {
        int t = a[i]; a[i] = a[j]; a[j] = t;
        i++; j--;
    }
}
\end{lstlisting}

\cost{$O(n)$ time, $O(1)$ space.}
\trap{Not normalising $k$. With \texttt{k = 7} on a length-5 array,
\texttt{reverse(a, 0, 6)} indexes past the end and throws.}

% ------------------------------------------------------------------
\problem{Merge two sorted arrays}{Easy}

Both inputs are sorted ascending. Return one sorted array with all the
elements.

\idea{One index into each; repeatedly take the smaller head, then drain
whichever array still has elements.}

\begin{dia}
  \node[note, anchor=east] at (-0.5,0.9) {a};
  \cellrowst{a}{0.9}{dead/1,soft/5,cell/9}
  \node[note, anchor=east] at (-0.5,0) {b};
  \cellrowst{b}{0}{dead/2,soft/3,cell/8}
  \node[ptr] at (0.8,1.5) {i};
  \node[ptr] at (0.8,-0.75) {j};
  \node[note, anchor=east] at (-0.5,-1.6) {out};
  \cellrowst{o}{-1.6}{hot/1,hot/2,cell/{},cell/{},cell/{},cell/{}}
  \node[note, anchor=west] at (4.2,-1.6) {next: compare 5 and 3};
\end{dia}

\textbf{How it works.} At every step the smallest element not yet taken must
be at the head of one of the two arrays, because each is sorted. So one
comparison is enough to decide what goes next --- there is no searching.

The two drain loops after the main one are not optional. The main loop exits
as soon as \emph{either} array runs out, and whatever remains in the other is
already sorted and larger than everything written so far, so it can be copied
straight across.

Using \texttt{<=} rather than \texttt{<} in the comparison keeps the merge
stable: equal elements from \texttt{a} come out before those from \texttt{b}.

\solheading
\begin{lstlisting}[style=java]
public static int[] merge(int[] a, int[] b) {
    if (a == null) a = new int[0];
    if (b == null) b = new int[0];
    int[] out = new int[a.length + b.length];
    int i = 0, j = 0, k = 0;
    while (i < a.length && j < b.length)
        out[k++] = (a[i] <= b[j]) ? a[i++] : b[j++];
    while (i < a.length) out[k++] = a[i++];
    while (j < b.length) out[k++] = b[j++];
    return out;
}
\end{lstlisting}

\cost{$O(n+m)$ time and space.}
\trap{Forgetting the drain loops, which silently truncates the output.}

% ------------------------------------------------------------------
\problem{Product of the array except self}{Medium}

Position $i$ of the result holds the product of every element except
\texttt{a[i]}. No division allowed.

\idea{Two sweeps. Left to right, store the product of everything \emph{before}
$i$. Right to left, multiply in the product of everything \emph{after} $i$,
carried in a single variable.}

\begin{dia}
  \node[note, anchor=east] at (-0.5,0.9) {input};
  \cellrowst{a}{0.9}{cell/1,cell/2,cell/3,cell/4}
  \node[note, anchor=east] at (-0.5,0) {prefix};
  \cellrowst{p}{0}{soft/1,soft/1,soft/2,soft/6}
  \draw[flow] (-0.1,0.45) -- (2.5,0.45);
  \node[note, anchor=east] at (-0.5,-0.9) {suffix};
  \cellrowst{s}{-0.9}{warm/24,warm/12,warm/4,warm/1}
  \draw[flow] (2.5,-0.45) -- (-0.1,-0.45);
  \node[note, anchor=east] at (-0.5,-1.8) {product};
  \cellrowst{r}{-1.8}{hot/24,hot/12,hot/8,hot/6}
\end{dia}

\textbf{How it works.} The product of everything except position $i$ splits
cleanly into two halves: what is to its left and what is to its right. Each
half is a running product, computable in one sweep.

The implementation is tighter than the diagram suggests. The first loop writes
the prefix products directly into the output array; the second loop walks
backwards, multiplying each slot by a running suffix carried in one variable.
So only the output array is allocated, and the extra space is $O(1)$.

\solheading
\begin{lstlisting}[style=java]
public static int[] productExceptSelf(int[] a) {
    if (a == null || a.length == 0) return new int[0];
    int n = a.length;
    int[] out = new int[n];
    out[0] = 1;
    for (int i = 1; i < n; i++) out[i] = out[i - 1] * a[i - 1];
    int suffix = 1;
    for (int i = n - 1; i >= 0; i--) {
        out[i] *= suffix;
        suffix *= a[i];
    }
    return out;
}
\end{lstlisting}

\cost{$O(n)$ time, $O(1)$ extra space beyond the output.}
\trap{The division shortcut --- total product divided by \texttt{a[i]} ---
breaks on one zero and is undefined on two. That is exactly why the statement
forbids it.}

% ------------------------------------------------------------------
\problem{Container with most water}{Medium}

Each element is the height of a vertical line. Pick two so the rectangle
between them holds the most water: the shorter height times the distance.

\idea{Start at the widest pair and close in, always moving the \emph{shorter}
line.}

\begin{dia}
  \foreach \x/\h in {0/1.8, 0.8/8, 1.6/6, 2.4/2, 3.2/5, 4.0/4, 4.8/8, 5.6/3, 6.4/7} {
    \draw[edge, line width=2pt] (\x,0) -- (\x,\h*0.17);}
  \draw[edge] (-0.3,0) -- (6.7,0);
  \fill[dcool, opacity=0.65] (0.8,0) rectangle (4.8,1.36);
  \draw[flow] (0.8,1.75) -- (0.8,1.45);
  \draw[flow] (4.8,1.75) -- (4.8,1.45);
  \node[ptr] at (0.8,1.95) {i};
  \node[ptr] at (4.8,1.95) {j};
  \caplabel{2.8}{-0.45}{area = $\min(8,8)\times 5 = 40$};
\end{dia}

\textbf{How it works.} The area is limited by the shorter of the two lines, so
think about what happens when you move each one inward. Moving the
\emph{taller} line loses width and cannot raise the limiting height --- the
shorter line is still there, still capping the area. So no better rectangle
uses that pair of positions, and the move is strictly wasteful.

Moving the shorter line, by contrast, at least gives the height a chance to
increase. That asymmetry is what makes it safe to discard one candidate per
step and finish in a single pass.

\solheading
\begin{lstlisting}[style=java]
public static int maxArea(int[] height) {
    if (height == null || height.length < 2) return 0;
    int i = 0, j = height.length - 1, best = 0;
    while (i < j) {
        int h = Math.min(height[i], height[j]);
        best = Math.max(best, h * (j - i));
        if (height[i] < height[j]) i++; else j--;
    }
    return best;
}
\end{lstlisting}

\cost{$O(n)$ time, $O(1)$ space.}
\trap{Moving whichever pointer you feel like, or both at once. The proof that
nothing is missed depends on always discarding the shorter line.}
}{}
\IfFileExists{ch3-hashing.tex}{\chapterpage{Hashing and counting}{Eight problems. If a question asks
\emph{how many}, \emph{have I seen}, or \emph{group by}, the answer starts
with a \texttt{HashMap}. The prefix trick in P18 and P21 is the highest-value
pattern in this book.}

% ------------------------------------------------------------------
\problem{Count pairs summing to a target}{Easy}

Count the pairs of indices $i<j$ with \texttt{a[i] + a[j] == target}.

\idea{Walk once. For each element, the number of earlier elements that complete
the pair is already in the map. Count first, then insert.}

\begin{dia}
  \cellrowst{a}{0}{dead/1,dead/5,dead/7,hot/1}
  \idxrow{-0.62}{3}
  \node[ptr] at (2.4,0.78) {x = 1};
  \draw[flow] (2.4,0.62) -- (2.4,0.44);
  \node[note, anchor=west] at (3.3,0.35) {target 6, so look up $6-1=5$};
  \node[note, anchor=west] at (3.3,-0.2) {map says \texttt{5} seen once};
  \node[note, anchor=west] at (3.3,-0.75) {$\Rightarrow$ one new pair};
  \node[soft, minimum width=11mm] at (0.0,-1.8) {1:1};
  \node[soft, minimum width=11mm] at (1.2,-1.8) {5:1};
  \node[soft, minimum width=11mm] at (2.4,-1.8) {7:1};
  \node[note, anchor=east] at (-0.7,-1.8) {map before};
\end{dia}

\textbf{How it works.} Every pair has a later element and an earlier one. If
you always count from the point of view of the later element, each pair is
counted exactly once --- no division by two, no double counting.

That is why the order of the two statements matters. Looking up the
complement \emph{before} inserting the current value means the map holds only
elements strictly to the left. Insert first and an element can pair with
itself: a single \texttt{3} with target 6 would report a pair that does not
exist.

\solheading
\begin{lstlisting}[style=java]
public static long countPairs(int[] a, int target) {
    if (a == null) return 0;
    Map<Integer,Integer> seen = new HashMap<>();
    long pairs = 0;
    for (int x : a) {
        pairs += seen.getOrDefault(target - x, 0);
        seen.merge(x, 1, Integer::sum);
    }
    return pairs;
}
\end{lstlisting}

\cost{$O(n)$ time, $O(n)$ space.}
\trap{Returning \texttt{int}. With $200\,000$ identical elements the count is
about $2\times10^{10}$, which overflows.}

% ------------------------------------------------------------------
\problem{Longest subarray with equal counts}{Medium}

Return the length of the longest contiguous subarray holding as many even
numbers as odd ones.

\idea{Turn it into a running balance: $+1$ for odd, $-1$ for even. A subarray is
balanced exactly when the balance at its two ends is equal.}

\begin{dia}
  \draw[edge] (0,0) -- (5.4,0);
  \node[idx, anchor=east] at (-0.15,0) {0};
  \node[idx, anchor=east] at (-0.15,0.55) {1};
  \node[idx, anchor=east] at (-0.15,1.1) {2};
  \draw[edge, densely dotted] (0,0.55) -- (5.4,0.55);
  \draw[flow, draw=accent, line width=1pt]
    (0,0) -- (0.6,0.55) -- (1.2,0) -- (1.8,0.55) -- (2.4,0) -- (3.0,0.55);
  \foreach \x/\y in {0/0, 0.6/0.55, 1.2/0, 1.8/0.55, 2.4/0, 3.0/0.55} {
    \fill[accent] (\x,\y) circle (1.6pt);}
  \node[idx] at (0,-0.35) {$-1$};
  \node[idx] at (0.6,-0.35) {0};
  \node[idx] at (1.2,-0.35) {1};
  \node[idx] at (1.8,-0.35) {2};
  \node[idx] at (2.4,-0.35) {3};
  \node[idx] at (3.0,-0.35) {4};
  \node[note, anchor=west] at (3.4,0.55) {index};
  \draw[span] (0,-0.85) -- (2.4,-0.85);
  \caplabel{1.2}{-1.35}{same balance at both ends $\Rightarrow$ length 4};
  \node[note, anchor=west] at (3.4,0) {balance};
\end{dia}

\textbf{How it works.} Write $B_i$ for the balance after the first $i$
elements. The subarray from $j+1$ to $i$ has as many odds as evens exactly
when $B_i = B_j$, because the difference $B_i - B_j$ counts (odds $-$ evens)
inside that window.

So the problem becomes: for each balance value, how far apart are its two
furthest occurrences? Since the balance is scanned left to right, you only
ever need the \emph{earliest} index at which each value appeared --- any later
repeat measured against it gives the longest span. Hence \texttt{putIfAbsent}
semantics: write the index only when the value is new.

The seed \texttt{firstIndex.put(0, -1)} represents the empty prefix. Without
it, a subarray that is balanced from index 0 has no earlier index to measure
against and is missed entirely.

\solheading
\begin{lstlisting}[style=java]
public static int longestBalanced(int[] nums) {
    if (nums == null || nums.length == 0) return 0;
    Map<Integer,Integer> firstIndex = new HashMap<>();
    firstIndex.put(0, -1);                 // the empty prefix
    int balance = 0, best = 0;
    for (int i = 0; i < nums.length; i++) {
        balance += (nums[i] % 2 != 0) ? 1 : -1;
        Integer seen = firstIndex.get(balance);
        if (seen == null) firstIndex.put(balance, i);
        else best = Math.max(best, i - seen);
    }
    return best;
}
\end{lstlisting}

\cost{$O(n)$ time, $O(n)$ space.}
\trap{Three in one problem. Overwriting the stored index on a repeat measures
the shortest span. Omitting the \texttt{0 -> -1} seed loses answers starting
at index 0. And \texttt{nums[i] \% 2 == 1} is false for $-3$.}
\alsoknown{Task 4 of the assessment; and with 0 and 1 in place of even and
odd, the classic ``contiguous array'' question.}

% ------------------------------------------------------------------
\problem{Group anagrams}{Medium}

Group the words that are anagrams of one another.

\noindent\texttt{["eat","tea","tan","ate"]} gives
\texttt{[[eat,tea,ate],[tan]]}.

\idea{Give every word a canonical key that all its anagrams share. Sorting its
letters is the simplest such key.}

\begin{dia}
  \node[cell, minimum width=13mm] (w1) at (0,1.2) {eat};
  \node[cell, minimum width=13mm] (w2) at (0,0.5) {tea};
  \node[cell, minimum width=13mm] (w3) at (0,-0.2) {tan};
  \node[cell, minimum width=13mm] (w4) at (0,-0.9) {ate};
  \node[soft, minimum width=14mm] (k1) at (2.6,0.85) {aet};
  \node[soft, minimum width=14mm] (k2) at (2.6,-0.55) {ant};
  \node[note] at (2.6,1.5) {sorted key};
  \node[hot, minimum width=30mm] (g1) at (5.6,0.85) {eat, tea, ate};
  \node[hot, minimum width=30mm] (g2) at (5.6,-0.55) {tan};
  \draw[flow] (w1) -- (k1); \draw[flow] (w2) -- (k1); \draw[flow] (w4) -- (k1);
  \draw[flow] (w3) -- (k2);
  \draw[flow] (k1) -- (g1); \draw[flow] (k2) -- (g2);
\end{dia}

\textbf{How it works.} Two words are anagrams precisely when they have the
same multiset of letters. A multiset has no natural identity in Java, so you
manufacture one: sort the letters and the result is identical for all members
of a group and different for everything else. That string becomes the map
key, and \texttt{computeIfAbsent} does the grouping in one line.

A 26-slot count rendered as a string is a faster key --- it drops the
$\log k$ sorting factor --- but sorting is shorter to write and almost always
fast enough.

\solheading
\begin{lstlisting}[style=java]
public static List<List<String>> groupAnagrams(String[] words) {
    Map<String,List<String>> groups = new HashMap<>();
    for (String w : words) {
        char[] letters = w.toCharArray();
        Arrays.sort(letters);
        String key = new String(letters);
        groups.computeIfAbsent(key, k -> new ArrayList<>()).add(w);
    }
    return new ArrayList<>(groups.values());
}
\end{lstlisting}

\cost{$O(n\,k\log k)$ for $n$ words of length $k$.}
\trap{Using the word's \texttt{hashCode} or the sum of its characters as the
key. Both collide --- \texttt{"ad"} and \texttt{"bc"} have the same letter
sum --- and the bug only shows on inputs you did not try.}

% ------------------------------------------------------------------
\problem{Top k frequent elements}{Medium}

Return the $k$ most frequent values, most frequent first.

\idea{Count, then select. A min-heap capped at size $k$ keeps the selection at
$O(n\log k)$.}

\begin{dia}
  \node[note, anchor=east] at (-0.6,0) {counts};
  \node[soft, minimum width=12mm] at (0,0) {1:3};
  \node[soft, minimum width=12mm] at (1.3,0) {2:2};
  \node[soft, minimum width=12mm] at (2.6,0) {3:1};
  \node[soft, minimum width=12mm] at (3.9,0) {4:5};
  \node[note, anchor=east] at (-0.6,-1.4) {heap, $k=2$};
  \node[hot, minimum width=12mm] (h1) at (0,-1.4) {2:2};
  \node[hot, minimum width=12mm] (h2) at (1.3,-1.4) {4:5};
  \node[dead, minimum width=12mm] (h3) at (2.9,-1.4) {3:1};
  \node[note, anchor=west] at (3.6,-1.4) {evicted: smallest count};
  \draw[back] (h1) -- (h3);
  \caplabel{0.65}{-2.1}{the root of a min-heap is the \emph{least} frequent};
\end{dia}

\textbf{How it works.} Counting is the easy half. The selection is where the
choice lies: sorting all distinct values is $O(d\log d)$, while a heap capped
at $k$ is $O(d\log k)$ and, more importantly, needs only $k$ entries in
memory --- which is what makes it work on a stream.

The counter-intuitive part is that you keep the top $k$ with a \emph{min}-heap.
Push everything; whenever the heap exceeds $k$, discard its root, which is the
least frequent thing currently held. What survives is the $k$ most frequent.
Because the root is the smallest, draining the heap yields ascending
frequency, so the output array is filled backwards.

\solheading
\begin{lstlisting}[style=java]
public static int[] topKFrequent(int[] a, int k) {
    Map<Integer,Integer> freq = new HashMap<>();
    for (int x : a) freq.merge(x, 1, Integer::sum);

    PriorityQueue<Map.Entry<Integer,Integer>> heap =
        new PriorityQueue<>(Map.Entry.comparingByValue());   // min-heap
    for (Map.Entry<Integer,Integer> e : freq.entrySet()) {
        heap.offer(e);
        if (heap.size() > k) heap.poll();     // drop the least frequent
    }
    int[] out = new int[heap.size()];
    for (int i = out.length - 1; i >= 0; i--) out[i] = heap.poll().getKey();
    return out;
}
\end{lstlisting}

\cost{$O(n\log k)$ time, $O(n)$ space for the counts.}
\trap{Filling the output forwards. The heap drains smallest-first, so a
forward fill returns the $k$ most frequent in exactly the wrong order.}

% ------------------------------------------------------------------
\problem{Subarray sum equals k}{Medium}

Count the contiguous subarrays summing to exactly $k$. Values may be negative.

\idea{Prefix sums. A window sums to $k$ whenever an earlier prefix equals the
current prefix minus $k$.}

\begin{dia}
  \cellrowst{a}{0.9}{cell/1,cell/2,cell/3,cell/-2}
  \node[note, anchor=east] at (-0.5,0.9) {input};
  \cellrowst{p}{0}{soft/0,soft/1,soft/3,hot/6,soft/4}
  \node[note, anchor=east] at (-0.5,0) {prefix};
  \node[idx] at (-0.0,-0.55) {$\varnothing$};
  \draw[span] (0.4,-1.0) -- (2.0,-1.0);
  \caplabel{1.2}{-1.5}{$6-1=5$: the window summing to 5 is $2+3$};
  \draw[back] (2.4,0.45) to[out=150,in=30] (0.8,0.45);
\end{dia}

\textbf{How it works.} Let $P_i$ be the sum of the first $i$ elements. The sum
of the window $(j,i]$ is $P_i - P_j$. Asking for a window summing to $k$ is
therefore asking: how many earlier prefixes equal $P_i - k$? Keep a map from
prefix value to how many times it has occurred, and the answer accumulates in
one pass.

The map is seeded with $0 \mapsto 1$ for the empty prefix, which is what lets
a window starting at index 0 be counted.

It is worth knowing why a sliding window fails here. Windows work when
growing the window can only increase the sum, so you can shrink from the left
when you overshoot. Negative values destroy that monotonicity; prefix sums do
not care.

\solheading
\begin{lstlisting}[style=java]
public static int subarraysSummingTo(int[] a, int k) {
    Map<Long,Integer> counts = new HashMap<>();
    counts.put(0L, 1);                     // the empty prefix
    long prefix = 0;
    int total = 0;
    for (int x : a) {
        prefix += x;
        total += counts.getOrDefault(prefix - k, 0);
        counts.merge(prefix, 1, Integer::sum);
    }
    return total;
}
\end{lstlisting}

\cost{$O(n)$ time, $O(n)$ space.}
\trap{Reaching for a sliding window because the problem says ``contiguous''.
Check for negative values first --- they are the tell.}

% ------------------------------------------------------------------
\problem{First duplicate}{Easy}

Return the first value that appears twice, scanning left to right, or $-1$.

\idea{A \texttt{HashSet}. \texttt{add} returns \texttt{false} when the element
was already there, so the check and the insert are a single call.}

\begin{dia}
  \cellrowst{a}{0}{dead/2,dead/1,dead/3,hot/1,cell/5}
  \idxrow{-0.62}{4}
  \node[soft, minimum width=9mm] at (0,-1.5) {2};
  \node[soft, minimum width=9mm] at (1.0,-1.5) {1};
  \node[soft, minimum width=9mm] at (2.0,-1.5) {3};
  \node[note, anchor=east] at (-0.5,-1.5) {set};
  \draw[flow] (2.4,-0.95) -- (1.1,-1.15);
  \caplabel{3.6}{-1.5}{\texttt{add(1)} returns false $\Rightarrow$ answer};
\end{dia}

\textbf{How it works.} \texttt{Set.add} is specified to return whether the set
changed. That single boolean replaces a \texttt{contains} followed by an
\texttt{add}, halving the number of hash lookups and removing the chance of
writing one without the other.

Be sure which question is being asked. This returns the value whose
\emph{second} occurrence comes earliest. ``The duplicate whose first
occurrence is earliest'' is a different answer and needs a second pass.

\solheading
\begin{lstlisting}[style=java]
public static int firstDuplicate(int[] a) {
    Set<Integer> seen = new HashSet<>();
    for (int x : a) if (!seen.add(x)) return x;
    return -1;
}
\end{lstlisting}

\cost{$O(n)$ time, $O(n)$ space.}
\trap{Misreading which ``first'' is wanted. Also, $-1$ as a sentinel is wrong
if $-1$ can itself be a duplicate value --- return \texttt{Integer} and
\texttt{null} if the statement allows negatives.}

% ------------------------------------------------------------------
\problem{Intersection of two arrays}{Easy}

Return the values present in both arrays, each once.

\idea{Put one array in a set, filter the other against it, and collect the hits
into a second set so duplicates collapse.}

\begin{dia}
  \draw[edge, fill=dsoft, fill opacity=0.6] (0,0) circle (1.1);
  \draw[edge, fill=dwarm, fill opacity=0.6] (1.5,0) circle (1.1);
  \node[note] at (-0.55,0) {2, 9};
  \node[note] at (0.75,0) {4, 5};
  \node[note] at (2.05,0) {7};
  \node[note] at (-0.55,1.3) {a};
  \node[note] at (2.05,1.3) {b};
  \node[note, anchor=west] at (3.2,0) {result: \texttt{[4, 5]}};
\end{dia}

\textbf{How it works.} The first set turns ``is this value in \texttt{a}'' from
a scan into an $O(1)$ question, which is the whole saving over the nested
loop. The second set exists for a different reason: if \texttt{b} contains a
shared value twice, a plain list would emit it twice.

\texttt{LinkedHashSet} is used rather than \texttt{HashSet} so the output
order is the order of first appearance in \texttt{b}. Plain \texttt{HashSet}
leaves the order unspecified, and graders usually compare arrays exactly.

\solheading
\begin{lstlisting}[style=java]
public static int[] intersection(int[] a, int[] b) {
    Set<Integer> inA = new HashSet<>();
    for (int x : a) inA.add(x);
    Set<Integer> both = new LinkedHashSet<>();
    for (int x : b) if (inA.contains(x)) both.add(x);
    int[] out = new int[both.size()];
    int i = 0;
    for (int x : both) out[i++] = x;
    return out;
}
\end{lstlisting}

\cost{$O(n+m)$ time, $O(n)$ space.}
\trap{Collecting into a \texttt{List}, which repeats values that recur in
\texttt{b}.}

% ------------------------------------------------------------------
\problem{Isomorphic strings}{Medium}

Two strings are isomorphic if the characters of one map onto the characters of
the other consistently and one-to-one. \texttt{"egg"} and \texttt{"add"} are;
\texttt{"foo"} and \texttt{"bar"} are not.

\idea{Two maps, one in each direction. One map alone permits two different
characters to land on the same target.}

\begin{dia}
  \node[cell] (e1) at (0,0.5) {f};
  \node[cell] (e2) at (0.8,0.5) {o};
  \node[cell] (e3) at (1.6,0.5) {o};
  \node[cell] (a1) at (0,-1.1) {b};
  \node[cell] (a2) at (0.8,-1.1) {a};
  \node[hot]  (a3) at (1.6,-1.1) {r};
  \draw[flow] (e1) -- (a1);
  \draw[flow] (e2) -- (a2);
  \draw[flow] (e3) -- (a3);
  \node[note, anchor=west] at (2.4,-0.3) {\texttt{o} already maps to \texttt{a},};
  \node[note, anchor=west] at (2.4,-0.8) {so \texttt{o}$\to$\texttt{r} contradicts it};
\end{dia}

\textbf{How it works.} Consistency in one direction says: the same source
character always produces the same target. One-to-one adds the reverse: no two
source characters share a target. \texttt{"ab"} and \texttt{"cc"} satisfies
the first and fails the second, which is why a single map is not enough.

\texttt{putIfAbsent} is doing double duty: it records the mapping if it is new
and returns the existing one if it is not, so the check and the write are one
call per direction.

\solheading
\begin{lstlisting}[style=java]
public static boolean isIsomorphic(String s, String t) {
    if (s == null || t == null || s.length() != t.length()) return false;
    Map<Character,Character> forward = new HashMap<>();
    Map<Character,Character> backward = new HashMap<>();
    for (int i = 0; i < s.length(); i++) {
        char a = s.charAt(i), b = t.charAt(i);
        Character f = forward.putIfAbsent(a, b);
        Character g = backward.putIfAbsent(b, a);
        if ((f != null && f != b) || (g != null && g != a)) return false;
    }
    return true;
}
\end{lstlisting}

\cost{$O(n)$ time, $O(1)$ space for a bounded alphabet.}
\trap{Comparing the boxed \texttt{Character} values with \texttt{!=} works here
only because the cache covers ASCII. With \texttt{Integer} keys above 127 the
identical code fails silently --- make unboxing a habit.}
}{}
\IfFileExists{ch4-sorting.tex}{\chapterpage{Sorting, intervals and comparators}{Seven problems. Writing a
correct multi-key comparator under time pressure is a skill in itself, and
interval sweeps are the most common ``medium'' in a screen.}

% ------------------------------------------------------------------
\problem{Sort by frequency, then by value}{Medium}

Sort the array so that more frequent values come first, and equally frequent
values come in ascending numeric order.
\texttt{[1,1,2,2,2,3]} gives \texttt{[2,2,2,1,1,3]}.

\idea{Count first, then sort the \emph{boxed} values with a two-key comparator.
Primitive arrays cannot take a comparator, so boxing is required.}

\solheading
\begin{lstlisting}[style=java]
public static int[] sortByFrequency(int[] a) {
    Map<Integer,Integer> freq = new HashMap<>();
    for (int x : a) freq.merge(x, 1, Integer::sum);

    Integer[] boxed = new Integer[a.length];
    for (int i = 0; i < a.length; i++) boxed[i] = a[i];

    Arrays.sort(boxed, (x, y) -> {
        int fx = freq.get(x), fy = freq.get(y);
        if (fx != fy) return Integer.compare(fy, fx);   // frequency desc
        return Integer.compare(x, y);                   // value asc
    });

    int[] out = new int[a.length];
    for (int i = 0; i < a.length; i++) out[i] = boxed[i];
    return out;
}
\end{lstlisting}

\cost{$O(n\log n)$ time, $O(n)$ space.}
\trap{\texttt{Arrays.sort(int[], Comparator)} does not exist --- there is no
such overload, and the compile error surprises people under pressure. Box to
\texttt{Integer[]} or sort a \texttt{List<Integer>}.}

% ------------------------------------------------------------------
\problem{Merge overlapping intervals}{Medium}

Given intervals as \texttt{[start, end]}, merge every pair that overlaps.
\texttt{[[1,3],[2,6],[8,10]]} gives \texttt{[[1,6],[8,10]]}.

\idea{Sort by start. Then one pass: either the next interval starts inside the
one you are building, in which case stretch its end, or it does not, in which
case close the current one and open a new one.}

\solheading
\begin{lstlisting}[style=java]
public static int[][] merge(int[][] intervals) {
    if (intervals == null || intervals.length == 0) return new int[0][];
    int[][] a = intervals.clone();
    Arrays.sort(a, (x, y) -> Integer.compare(x[0], y[0]));

    List<int[]> out = new ArrayList<>();
    int[] current = new int[] {a[0][0], a[0][1]};
    for (int i = 1; i < a.length; i++) {
        if (a[i][0] <= current[1]) {
            current[1] = Math.max(current[1], a[i][1]);
        } else {
            out.add(current);
            current = new int[] {a[i][0], a[i][1]};
        }
    }
    out.add(current);
    return out.toArray(new int[0][]);
}
\end{lstlisting}

\cost{$O(n\log n)$ time, $O(n)$ space.}
\trap{\texttt{current[1] = a[i][1]} without the \texttt{Math.max} loses a
fully contained interval: merging \texttt{[1,10]} with \texttt{[2,3]} must
still end at 10. Also remember to add the last \texttt{current} after the
loop.}

% ------------------------------------------------------------------
\problem{Can one person attend all meetings?}{Easy}

Given meeting intervals, return \texttt{true} if none of them overlap.

\idea{Sort by start; any overlap is then between neighbours.}

\solheading
\begin{lstlisting}[style=java]
public static boolean canAttendAll(int[][] meetings) {
    if (meetings == null || meetings.length < 2) return true;
    int[][] a = meetings.clone();
    Arrays.sort(a, (x, y) -> Integer.compare(x[0], y[0]));
    for (int i = 1; i < a.length; i++)
        if (a[i][0] < a[i - 1][1]) return false;
    return true;
}
\end{lstlisting}

\cost{$O(n\log n)$ time.}
\trap{Whether touching intervals count as overlapping. \texttt{[1,5]} and
\texttt{[5,9]} do not overlap if the end is exclusive --- hence \texttt{<} and
not \texttt{<=}. The statement decides; read it.}

% ------------------------------------------------------------------
\problem{Minimum number of meeting rooms}{Medium}

Return the largest number of meetings running at the same time.

\idea{A sweep. Put every start and every end on a timeline, sort both, then walk
them: a start raises the count, an end lowers it. Track the peak. Ends must
be processed before starts at the same instant, since a room freed at 5 can
be reused at 5.}

\solheading
\begin{lstlisting}[style=java]
public static int minRooms(int[][] meetings) {
    if (meetings == null || meetings.length == 0) return 0;
    int n = meetings.length;
    int[] starts = new int[n], ends = new int[n];
    for (int i = 0; i < n; i++) {
        starts[i] = meetings[i][0];
        ends[i]   = meetings[i][1];
    }
    Arrays.sort(starts);
    Arrays.sort(ends);

    int rooms = 0, peak = 0, e = 0;
    for (int s = 0; s < n; s++) {
        while (e < n && ends[e] <= starts[s]) { rooms--; e++; }
        rooms++;
        peak = Math.max(peak, rooms);
    }
    return peak;
}
\end{lstlisting}

\cost{$O(n\log n)$ time, $O(n)$ space.}
\trap{\texttt{ends[e] < starts[s]} instead of \texttt{<=} books a second room
for back-to-back meetings. The two arrays are sorted independently on
purpose --- the pairing between a particular start and end does not matter.}

% ------------------------------------------------------------------
\problem{Insert an interval}{Medium}

The intervals are already sorted and non-overlapping. Insert a new one, merging
where it overlaps.

\idea{Three phases, no sorting needed: copy everything that ends before the new
one starts, absorb everything that overlaps into a widened interval, then
copy the rest.}

\solheading
\begin{lstlisting}[style=java]
public static int[][] insert(int[][] intervals, int[] newInterval) {
    List<int[]> out = new ArrayList<>();
    int i = 0, n = intervals.length;
    int start = newInterval[0], end = newInterval[1];

    while (i < n && intervals[i][1] < start) out.add(intervals[i++]);

    while (i < n && intervals[i][0] <= end) {
        start = Math.min(start, intervals[i][0]);
        end   = Math.max(end,   intervals[i][1]);
        i++;
    }
    out.add(new int[] {start, end});

    while (i < n) out.add(intervals[i++]);
    return out.toArray(new int[0][]);
}
\end{lstlisting}

\cost{$O(n)$ time --- the input is already sorted, so do not re-sort it.}
\trap{Adding the merged interval inside the second loop, which emits it once
per overlap. It goes in exactly once, after the loop.}

% ------------------------------------------------------------------
\problem{Sort employees by salary then name}{Easy}

Sort records by salary descending, and alphabetically for equal salaries. A
pure comparator drill --- this exact shape appears in data-processing tasks.

\idea{The \texttt{Comparator} factory methods compose. Build the ordering
declaratively rather than hand-rolling the branches.}

\solheading
\begin{lstlisting}[style=java]
record Employee(String name, int salary) {}

public static void sortEmployees(List<Employee> staff) {
    staff.sort(Comparator.comparingInt(Employee::salary).reversed()
                         .thenComparing(Employee::name));
}

// the hand-written equivalent, when you cannot remember the fluent form
public static void sortEmployeesManually(List<Employee> staff) {
    staff.sort((a, b) -> {
        if (a.salary() != b.salary())
            return Integer.compare(b.salary(), a.salary());
        return a.name().compareTo(b.name());
    });
}
\end{lstlisting}

\cost{$O(n\log n)$.}
\trap{\texttt{.reversed()} reverses \emph{everything composed so far}. Chaining
it on after \texttt{thenComparing(name)} reverses the name order as well. Put
\texttt{reversed()} on the salary key alone, as above.}

% ------------------------------------------------------------------
\problem{Kth largest element}{Medium}

Return the $k$th largest value in an unsorted array (the $k$th in sorted
descending order, duplicates counted).

\idea{Keep a min-heap of size $k$. The smallest thing in it is, once you have
seen everything, the $k$th largest overall.}

\solheading
\begin{lstlisting}[style=java]
public static int kthLargest(int[] a, int k) {
    PriorityQueue<Integer> heap = new PriorityQueue<>();   // min-heap
    for (int x : a) {
        heap.offer(x);
        if (heap.size() > k) heap.poll();
    }
    return heap.peek();
}
\end{lstlisting}

\cost{$O(n\log k)$ time, $O(k)$ space. Sorting is $O(n\log n)$ and perfectly
acceptable when $k$ is close to $n$.}
\trap{Using a max-heap and polling $k$ times is $O(n + k\log n)$ to build and
drain --- fine, but it needs all $n$ elements in memory. The size-$k$ min-heap
is the answer when the input is a stream.}
}{}
\IfFileExists{ch5-matrix.tex}{\chapterpage{Matrices}{Six problems. Two-dimensional arrays are mostly an
index-discipline test: get \texttt{rows} and \texttt{cols} the right way round
and most of the difficulty disappears.}

% ------------------------------------------------------------------
\problem{Rotate a square matrix 90 degrees clockwise}{Medium}

Rotate in place. \texttt{[[1,2],[3,4]]} becomes \texttt{[[3,1],[4,2]]}.

\idea{Transpose (reflect across the main diagonal), then reverse each row. Two
simple passes beat one confusing index formula.}

\solheading
\begin{lstlisting}[style=java]
public static void rotate90(int[][] m) {
    int n = m.length;
    for (int i = 0; i < n; i++)
        for (int j = i + 1; j < n; j++) {      // j starts at i+1
            int t = m[i][j]; m[i][j] = m[j][i]; m[j][i] = t;
        }
    for (int[] row : m)
        for (int i = 0, j = n - 1; i < j; i++, j--) {
            int t = row[i]; row[i] = row[j]; row[j] = t;
        }
}
\end{lstlisting}

\cost{$O(n^2)$ time, $O(1)$ space.}
\trap{Starting the inner transpose loop at \texttt{j = 0} swaps every pair
twice and leaves the matrix unchanged. It must start at \texttt{i + 1}.}

% ------------------------------------------------------------------
\problem{Spiral traversal}{Medium}

Return the elements in spiral order: left to right along the top, down the
right side, right to left along the bottom, up the left side, then inwards.

\idea{Four moving boundaries. After walking each edge, shrink that boundary, and
check before the two ``return'' edges that the boundaries have not crossed.}

\solheading
\begin{lstlisting}[style=java]
public static List<Integer> spiral(int[][] m) {
    List<Integer> out = new ArrayList<>();
    if (m == null || m.length == 0) return out;
    int top = 0, bottom = m.length - 1;
    int left = 0, right = m[0].length - 1;

    while (top <= bottom && left <= right) {
        for (int c = left; c <= right; c++) out.add(m[top][c]);
        top++;
        for (int r = top; r <= bottom; r++) out.add(m[r][right]);
        right--;
        if (top <= bottom) {
            for (int c = right; c >= left; c--) out.add(m[bottom][c]);
            bottom--;
        }
        if (left <= right) {
            for (int r = bottom; r >= top; r--) out.add(m[r][left]);
            left++;
        }
    }
    return out;
}
\end{lstlisting}

\cost{$O(\text{rows}\times\text{cols})$ time.}
\trap{Omitting the two guards duplicates a row or a column on non-square
matrices --- a single remaining row would be walked both left-to-right and
right-to-left.}

% ------------------------------------------------------------------
\problem{Count islands}{Medium}

The grid holds \texttt{'1'} for land and \texttt{'0'} for water. Count the
connected groups of land, using four-way adjacency.

\idea{Scan every cell. The first time you touch a piece of land, increment the
count and flood the whole island so it is never counted again.}

\solheading
\begin{lstlisting}[style=java]
public static int countIslands(char[][] grid) {
    if (grid == null || grid.length == 0) return 0;
    int rows = grid.length, cols = grid[0].length, islands = 0;
    int[][] dirs = {{-1,0},{1,0},{0,-1},{0,1}};

    for (int r = 0; r < rows; r++)
        for (int c = 0; c < cols; c++) {
            if (grid[r][c] != '1') continue;
            islands++;
            Deque<int[]> queue = new ArrayDeque<>();
            queue.add(new int[] {r, c});
            grid[r][c] = '0';                   // mark on push
            while (!queue.isEmpty()) {
                int[] cell = queue.poll();
                for (int[] d : dirs) {
                    int nr = cell[0] + d[0], nc = cell[1] + d[1];
                    if (nr < 0 || nr >= rows || nc < 0 || nc >= cols) continue;
                    if (grid[nr][nc] != '1') continue;
                    grid[nr][nc] = '0';
                    queue.add(new int[] {nr, nc});
                }
            }
        }
    return islands;
}
\end{lstlisting}

\cost{$O(\text{rows}\times\text{cols})$ time and space.}
\trap{Marking a cell visited when you \emph{pop} it rather than when you
\emph{push} it lets the same cell enter the queue many times, and a large
open grid then blows up. Mark on push, always.}

% ------------------------------------------------------------------
\problem{Set matrix zeroes}{Medium}

If any cell is zero, set its whole row and column to zero. The catch is doing
it without a second matrix.

\idea{Two passes. First record \emph{which} rows and columns must be cleared;
only then clear them. Clearing as you go turns the whole matrix to zero.}

\solheading
\begin{lstlisting}[style=java]
public static void setZeroes(int[][] m) {
    if (m == null || m.length == 0) return;
    int rows = m.length, cols = m[0].length;
    boolean[] clearRow = new boolean[rows];
    boolean[] clearCol = new boolean[cols];

    for (int r = 0; r < rows; r++)
        for (int c = 0; c < cols; c++)
            if (m[r][c] == 0) { clearRow[r] = true; clearCol[c] = true; }

    for (int r = 0; r < rows; r++)
        for (int c = 0; c < cols; c++)
            if (clearRow[r] || clearCol[c]) m[r][c] = 0;
}
\end{lstlisting}

\cost{$O(\text{rows}\times\text{cols})$ time, $O(\text{rows}+\text{cols})$
space.}
\trap{The naive one-pass version writes zeros that the rest of the pass then
reads as original zeros, cascading until everything is zero. The $O(1)$-space
variant stores the flags in row 0 and column 0 and needs a separate flag for
each of those --- more clever than it is worth under time pressure.}

% ------------------------------------------------------------------
\problem{Diagonal sums}{Easy}

Sum the main diagonal and the anti-diagonal of a square matrix, counting the
centre of an odd-sized matrix once.

\idea{One loop. At step $i$ the two diagonal cells are $m[i][i]$ and
$m[i][n-1-i]$; they coincide exactly when $i = n-1-i$.}

\solheading
\begin{lstlisting}[style=java]
public static long diagonalSum(int[][] m) {
    int n = m.length;
    long total = 0;
    for (int i = 0; i < n; i++) {
        total += m[i][i];
        if (i != n - 1 - i) total += m[i][n - 1 - i];
    }
    return total;
}
\end{lstlisting}

\cost{$O(n)$ time, $O(1)$ space.}
\trap{Double-counting the centre cell of an odd matrix. And the total can
exceed \texttt{int} --- use \texttt{long}.}

% ------------------------------------------------------------------
\problem{Search a row-and-column sorted matrix}{Medium}

Every row is sorted left to right and every column top to bottom. Decide
whether a target is present.

\idea{Start at the top-right corner. That cell is the largest in its row and the
smallest in its column, so a comparison always eliminates an entire row or an
entire column.}

\solheading
\begin{lstlisting}[style=java]
public static boolean contains(int[][] m, int target) {
    if (m == null || m.length == 0) return false;
    int r = 0, c = m[0].length - 1;
    while (r < m.length && c >= 0) {
        if (m[r][c] == target) return true;
        if (m[r][c] > target) c--;     // whole column is too big
        else r++;                      // whole row is too small
    }
    return false;
}
\end{lstlisting}

\cost{$O(\text{rows}+\text{cols})$ time, $O(1)$ space.}
\trap{Starting at the top-left corner gives no information: both moves lead to
larger values, so you cannot eliminate anything. The corner choice is the
entire trick.}
}{}
\IfFileExists{ch6-recursion.tex}{\chapterpage{Recursion, BFS and DFS}{Seven problems. Recursion is fine for
trees, where depth is logarithmic. On grids, use an explicit queue --- depth
there is measured in tens of thousands.}

% ------------------------------------------------------------------
\problem{Fibonacci with memoisation}{Easy}

Return the $n$th Fibonacci number. The point is the memo, not the formula.

\idea{The naive recursion recomputes the same subproblems exponentially often.
Cache each answer the first time you compute it. Better still, note that only
two previous values are ever needed and drop the recursion entirely.}

\solheading
\begin{lstlisting}[style=java]
public static long fib(int n) {
    if (n < 2) return n;
    long prev = 0, current = 1;
    for (int i = 2; i <= n; i++) {
        long next = prev + current;
        prev = current;
        current = next;
    }
    return current;
}

// the memoised recursion, when the recurrence is not so easily unrolled
public static long fibMemo(int n, Map<Integer,Long> memo) {
    if (n < 2) return n;
    Long cached = memo.get(n);
    if (cached != null) return cached;
    long value = fibMemo(n - 1, memo) + fibMemo(n - 2, memo);
    memo.put(n, value);
    return value;
}
\end{lstlisting}

\cost{Iterative: $O(n)$ time, $O(1)$ space. Memoised: $O(n)$ both.}
\trap{\texttt{int} overflows at \texttt{fib(47)}. If the statement does not
bound $n$, use \texttt{long} and say so.}

% ------------------------------------------------------------------
\problem{Maximum depth of a binary tree}{Easy}

Return the number of nodes on the longest root-to-leaf path.

\idea{The depth of a node is one more than the deeper of its two subtrees. This
is the template every tree recursion follows: base case for \texttt{null},
combine the children.}

\solheading
\begin{lstlisting}[style=java]
static class Node {
    int value;
    Node left, right;
}

public static int maxDepth(Node root) {
    if (root == null) return 0;
    return 1 + Math.max(maxDepth(root.left), maxDepth(root.right));
}
\end{lstlisting}

\cost{$O(n)$ time, $O(h)$ stack for height $h$.}
\trap{Returning \texttt{1} rather than \texttt{0} for a \texttt{null} node
counts a missing child as a level and gives every leaf a depth of 2.}

% ------------------------------------------------------------------
\problem{Level-order traversal}{Medium}

Return the node values grouped by depth, top level first.

\idea{BFS with a queue, taking a whole level at a time: record the queue size
before the loop, then pop exactly that many nodes. That snapshot is what
separates one level from the next.}

\solheading
\begin{lstlisting}[style=java]
public static List<List<Integer>> levelOrder(Node root) {
    List<List<Integer>> out = new ArrayList<>();
    if (root == null) return out;
    Deque<Node> queue = new ArrayDeque<>();
    queue.add(root);
    while (!queue.isEmpty()) {
        int levelSize = queue.size();            // snapshot
        List<Integer> level = new ArrayList<>();
        for (int i = 0; i < levelSize; i++) {
            Node node = queue.poll();
            level.add(node.value);
            if (node.left != null) queue.add(node.left);
            if (node.right != null) queue.add(node.right);
        }
        out.add(level);
    }
    return out;
}
\end{lstlisting}

\cost{$O(n)$ time, $O(w)$ space for the widest level $w$.}
\trap{Writing \texttt{for (int i = 0; i < queue.size(); i++)} re-evaluates the
size as children are added, so levels bleed into each other. Take the
snapshot into a variable.}

% ------------------------------------------------------------------
\problem{Flood fill}{Easy}

Starting at a pixel, repaint it and every same-coloured pixel connected to it
(four-way) with a new colour.

\idea{BFS from the start pixel. Guard against the new colour equalling the old
one, or the fill never terminates.}

\solheading
\begin{lstlisting}[style=java]
public static int[][] floodFill(int[][] image, int sr, int sc, int colour) {
    int old = image[sr][sc];
    if (old == colour) return image;             // the essential guard
    int rows = image.length, cols = image[0].length;
    int[][] dirs = {{-1,0},{1,0},{0,-1},{0,1}};

    Deque<int[]> queue = new ArrayDeque<>();
    queue.add(new int[] {sr, sc});
    image[sr][sc] = colour;
    while (!queue.isEmpty()) {
        int[] cell = queue.poll();
        for (int[] d : dirs) {
            int r = cell[0] + d[0], c = cell[1] + d[1];
            if (r < 0 || r >= rows || c < 0 || c >= cols) continue;
            if (image[r][c] != old) continue;
            image[r][c] = colour;
            queue.add(new int[] {r, c});
        }
    }
    return image;
}
\end{lstlisting}

\cost{$O(\text{rows}\times\text{cols})$.}
\trap{Without the \texttt{old == colour} early return, every repainted cell
still matches the target colour, so it is enqueued again forever.}
\alsoknown{The engine behind Task 3 of the assessment.}

% ------------------------------------------------------------------
\problem{All permutations of a string}{Medium}

Return every ordering of the characters. For $n$ characters there are $n!$ of
them, so this is only ever asked for small inputs.

\idea{Backtracking. Choose an unused character, recurse, then put it back. The
``put it back'' line is what makes it explore every branch rather than one.}

\solheading
\begin{lstlisting}[style=java]
public static List<String> permutations(String s) {
    List<String> out = new ArrayList<>();
    if (s == null || s.isEmpty()) return out;
    build(s.toCharArray(), new boolean[s.length()],
          new StringBuilder(), out);
    return out;
}

private static void build(char[] chars, boolean[] used,
                          StringBuilder current, List<String> out) {
    if (current.length() == chars.length) {
        out.add(current.toString());
        return;
    }
    for (int i = 0; i < chars.length; i++) {
        if (used[i]) continue;
        used[i] = true;
        current.append(chars[i]);
        build(chars, used, current, out);
        current.deleteCharAt(current.length() - 1);   // undo
        used[i] = false;                              // undo
    }
}
\end{lstlisting}

\cost{$O(n \cdot n!)$ time.}
\trap{Forgetting either undo. The \texttt{StringBuilder} and the
\texttt{used} array are shared across the whole search, so every change must
be reversed on the way out.}

% ------------------------------------------------------------------
\problem{All subsets}{Medium}

Return every subset of a set of distinct numbers --- $2^n$ of them.

\idea{Each element is either in or out, so the subsets map one-to-one onto the
binary numbers from $0$ to $2^n - 1$. Bit $i$ set means element $i$ is
included. No recursion needed.}

\solheading
\begin{lstlisting}[style=java]
public static List<List<Integer>> subsets(int[] a) {
    List<List<Integer>> out = new ArrayList<>();
    int total = 1 << a.length;                   // 2^n
    for (int mask = 0; mask < total; mask++) {
        List<Integer> subset = new ArrayList<>();
        for (int i = 0; i < a.length; i++)
            if ((mask & (1 << i)) != 0) subset.add(a[i]);
        out.add(subset);
    }
    return out;
}
\end{lstlisting}

\cost{$O(n \cdot 2^n)$ time.}
\trap{\texttt{1 << a.length} overflows once the array has 31 elements. That
is academic here --- $2^{31}$ subsets would not fit in memory anyway --- but
the shift is the thing that breaks first.}

% ------------------------------------------------------------------
\problem{Shortest path in a grid with obstacles}{Medium}

Cells are open (\texttt{0}) or blocked (\texttt{1}). Return the fewest steps
from the top-left to the bottom-right, moving four-way, or $-1$.

\idea{BFS, not DFS. On an unweighted graph BFS reaches every cell by a shortest
path, so the first time you arrive at the target you are done.}

\solheading
\begin{lstlisting}[style=java]
public static int shortestPath(int[][] grid) {
    if (grid == null || grid.length == 0) return -1;
    int rows = grid.length, cols = grid[0].length;
    if (grid[0][0] == 1 || grid[rows - 1][cols - 1] == 1) return -1;
    int[][] dirs = {{-1,0},{1,0},{0,-1},{0,1}};

    boolean[][] seen = new boolean[rows][cols];
    Deque<int[]> queue = new ArrayDeque<>();
    queue.add(new int[] {0, 0});
    seen[0][0] = true;
    int steps = 1;

    while (!queue.isEmpty()) {
        int levelSize = queue.size();
        for (int i = 0; i < levelSize; i++) {
            int[] cell = queue.poll();
            if (cell[0] == rows - 1 && cell[1] == cols - 1) return steps;
            for (int[] d : dirs) {
                int r = cell[0] + d[0], c = cell[1] + d[1];
                if (r < 0 || r >= rows || c < 0 || c >= cols) continue;
                if (seen[r][c] || grid[r][c] == 1) continue;
                seen[r][c] = true;
                queue.add(new int[] {r, c});
            }
        }
        steps++;
    }
    return -1;
}
\end{lstlisting}

\cost{$O(\text{rows}\times\text{cols})$ time and space.}
\trap{DFS finds \emph{a} path, not the shortest one, and people reach for it
because recursion is quicker to type. If the question says shortest and the
edges are unweighted, it is BFS.}
}{}
\IfFileExists{ch7-stacks.tex}{\chapterpage{Stacks and queues}{Five problems. A stack is the right answer
whenever the most recent unresolved thing is the one you need next.}

% ------------------------------------------------------------------
\problem{Valid parentheses}{Easy}

Given a string of \texttt{()}, \texttt{[]} and \texttt{\{\}}, decide whether
every bracket is closed by the matching kind, in the right order.

\idea{Push openers. On a closer, the top of the stack must be its partner. At
the end the stack must be empty.}

\solheading
\begin{lstlisting}[style=java]
public static boolean isValid(String s) {
    if (s == null) return true;
    Deque<Character> stack = new ArrayDeque<>();
    for (char c : s.toCharArray()) {
        switch (c) {
            case '(' : case '[' : case '{' : stack.push(c); break;
            case ')' : if (stack.isEmpty() || stack.pop() != '(') return false; break;
            case ']' : if (stack.isEmpty() || stack.pop() != '[') return false; break;
            case '}' : if (stack.isEmpty() || stack.pop() != '{') return false; break;
            default  : break;    // ignore anything else
        }
    }
    return stack.isEmpty();
}
\end{lstlisting}

\cost{$O(n)$ time, $O(n)$ space.}
\trap{Two endings get forgotten. Popping an empty stack on a stray closer
throws, so test \texttt{isEmpty()} first; and leaving the stack non-empty at
the end means unclosed openers, which is also invalid.}

% ------------------------------------------------------------------
\problem{Evaluate reverse Polish notation}{Medium}

Evaluate tokens like \texttt{["2","1","+","3","*"]}, which means
$(2+1)\times 3 = 9$.

\idea{Push numbers. On an operator, pop two, apply, push the result. The first
value popped is the right-hand operand.}

\solheading
\begin{lstlisting}[style=java]
public static int evalRPN(String[] tokens) {
    Deque<Integer> stack = new ArrayDeque<>();
    for (String t : tokens) {
        switch (t) {
            case "+": case "-": case "*": case "/": {
                int right = stack.pop();           // order matters
                int left  = stack.pop();
                stack.push(apply(left, right, t));
                break;
            }
            default:
                stack.push(Integer.parseInt(t));
        }
    }
    return stack.pop();
}

private static int apply(int a, int b, String op) {
    switch (op) {
        case "+": return a + b;
        case "-": return a - b;
        case "*": return a * b;
        default : return a / b;
    }
}
\end{lstlisting}

\cost{$O(n)$ time, $O(n)$ space.}
\trap{Popping the operands in the wrong order. It is invisible for
\texttt{+} and \texttt{*} and wrong for \texttt{-} and \texttt{/}, so the bug
survives your first two test cases.}

% ------------------------------------------------------------------
\problem{Next greater element}{Medium}

For every element, find the first element to its right that is larger, or
$-1$. \texttt{[2,1,2,4]} gives \texttt{[4,2,4,-1]}.

\idea{A monotonic stack of \emph{indices} still waiting for an answer. When a
new value arrives, it resolves every waiting index smaller than it. Each
index is pushed and popped once.}

\solheading
\begin{lstlisting}[style=java]
public static int[] nextGreater(int[] a) {
    int n = a.length;
    int[] out = new int[n];
    Arrays.fill(out, -1);
    Deque<Integer> waiting = new ArrayDeque<>();   // indices

    for (int i = 0; i < n; i++) {
        while (!waiting.isEmpty() && a[waiting.peek()] < a[i])
            out[waiting.pop()] = a[i];
        waiting.push(i);
    }
    return out;
}
\end{lstlisting}

\cost{$O(n)$ time --- the inner \texttt{while} is amortised, since every index
leaves the stack at most once. $O(n)$ space.}
\trap{The nested-loop version is $O(n^2)$ and looks so natural that people
write it without noticing. Also store indices, not values, or you cannot
write the answer back to the right slot.}
\alsoknown{``Daily temperatures'', where the answer is the gap \texttt{i - j}
rather than the value.}

% ------------------------------------------------------------------
\problem{Min stack}{Medium}

A stack that also reports its minimum in $O(1)$.

\idea{Keep a second stack holding the minimum \emph{as of} each push. Pushing
records the smaller of the new value and the current minimum; popping
discards both in step, so the two stacks stay aligned.}

\solheading
\begin{lstlisting}[style=java]
class MinStack {
    private final Deque<Integer> values = new ArrayDeque<>();
    private final Deque<Integer> minima = new ArrayDeque<>();

    void push(int x) {
        values.push(x);
        minima.push(minima.isEmpty() ? x : Math.min(x, minima.peek()));
    }

    int pop() {
        minima.pop();
        return values.pop();
    }

    int top()     { return values.peek(); }
    int getMin()  { return minima.peek(); }
}
\end{lstlisting}

\cost{$O(1)$ for every operation, $O(n)$ space.}
\trap{Keeping a single \texttt{int min} field instead. It is right until you
pop the minimum, and then there is no way to recover the previous one without
scanning.}

% ------------------------------------------------------------------
\problem{Simplify a Unix path}{Medium}

Turn \texttt{"/a/./b/../../c/"} into \texttt{"/c"}. Handle \texttt{.},
\texttt{..} and repeated slashes.

\idea{Split on slashes and push each real component. \texttt{..} pops,
\texttt{.} and empty components are skipped. Join what is left.}

\solheading
\begin{lstlisting}[style=java]
public static String simplify(String path) {
    Deque<String> parts = new ArrayDeque<>();     // used as a list
    for (String part : path.split("/")) {
        if (part.isEmpty() || part.equals(".")) continue;
        if (part.equals("..")) parts.pollLast();  // no-op when empty
        else parts.addLast(part);
    }
    if (parts.isEmpty()) return "/";
    StringBuilder out = new StringBuilder();
    for (String part : parts) out.append('/').append(part);
    return out.toString();
}
\end{lstlisting}

\cost{$O(n)$ time, $O(n)$ space.}
\trap{\texttt{..} at the root must do nothing, not throw ---
\texttt{pollLast} returns \texttt{null} on an empty deque where
\texttt{removeLast} throws. And the result for an empty stack is
\texttt{"/"}, not \texttt{""}.}
}{}
\IfFileExists{ch8-data.tex}{\chapterpage{Data processing}{Six problems in the shape candidates actually
report from consulting-tech screens: parse records, validate them, group them,
aggregate them, report the top few. Less algorithmic than the rest of the book
and more likely to be what you face.}

% ------------------------------------------------------------------
\problem{Average salary by department}{Medium}

Records are \texttt{"name,department,salary"}. Return the top $n$ departments
by average salary as \texttt{"department(average)"}, average descending then
name ascending. The average is integer division. Ignore any record that does
not have three fields, has an empty department, or whose salary is not a whole
number at least zero.

\idea{Accumulate a \emph{total and a count} per department, never a running
average. Compute each average once, at the end, from exact numbers.}

\solheading
\begin{lstlisting}[style=java]
public static String[] topDepartments(String[] records, int n) {
    if (records == null || n <= 0) return new String[0];
    Map<String,long[]> tally = new HashMap<>();      // [total, count]

    for (String record : records) {
        if (record == null) continue;
        String[] parts = record.split(",", -1);      // keep empty fields
        if (parts.length != 3) continue;
        String dept = parts[1].trim();
        if (dept.isEmpty()) continue;
        long salary;
        try { salary = Long.parseLong(parts[2].trim()); }
        catch (NumberFormatException e) { continue; }
        if (salary < 0) continue;

        long[] slot = tally.computeIfAbsent(dept, k -> new long[2]);
        slot[0] += salary;
        slot[1]++;
    }
    if (tally.isEmpty()) return new String[0];

    List<String> names = new ArrayList<>(tally.keySet());
    names.sort((a, b) -> {
        long avgA = tally.get(a)[0] / tally.get(a)[1];
        long avgB = tally.get(b)[0] / tally.get(b)[1];
        if (avgA != avgB) return Long.compare(avgB, avgA);
        return a.compareTo(b);
    });

    int size = Math.min(n, names.size());
    String[] out = new String[size];
    for (int i = 0; i < size; i++) {
        long[] slot = tally.get(names.get(i));
        out[i] = names.get(i) + "(" + (slot[0] / slot[1]) + ")";
    }
    return out;
}
\end{lstlisting}

\cost{$O(R + D\log D)$ for $R$ records and $D$ departments.}
\trap{Four of them, and each one is a real submission killer.
\texttt{split(",")} without \texttt{-1} silently drops trailing empty fields,
so \texttt{"a,x,"} looks like a two-field record instead of an invalid
three-field one. An \texttt{int} total overflows. Averaging averages is wrong
whenever departments differ in size. And a negative salary parses perfectly
well, so \texttt{parseInt} succeeding is not the same as the record being
valid.}
\alsoknown{Task 2 of the assessment.}

% ------------------------------------------------------------------
\problem{Filter and total transactions}{Medium}

Records are \texttt{"id,category,amount"} with the amount in cents. Return the
total per category, but only for categories whose total is at least a
threshold, sorted by total descending.

\idea{One grouping pass, then filter, then sort. Doing the filter before the
grouping is finished is the common error --- a category only crosses the
threshold once all its rows are in.}

\solheading
\begin{lstlisting}[style=java]
public static List<String> categoriesOver(String[] rows, long threshold) {
    Map<String,Long> totals = new HashMap<>();
    for (String row : rows) {
        if (row == null) continue;
        String[] parts = row.split(",", -1);
        if (parts.length != 3) continue;
        try {
            totals.merge(parts[1].trim(),
                         Long.parseLong(parts[2].trim()), Long::sum);
        } catch (NumberFormatException ignored) { }
    }

    List<Map.Entry<String,Long>> kept = new ArrayList<>();
    for (Map.Entry<String,Long> e : totals.entrySet())
        if (e.getValue() >= threshold) kept.add(e);

    kept.sort((a, b) -> b.getValue().compareTo(a.getValue()));

    List<String> out = new ArrayList<>();
    for (Map.Entry<String,Long> e : kept)
        out.add(e.getKey() + "=" + e.getValue());
    return out;
}
\end{lstlisting}

\cost{$O(R + C\log C)$.}
\trap{Filtering inside the parsing loop. Also, \texttt{b.getValue().compareTo}
on boxed \texttt{Long} is fine, but \texttt{(int)(b.getValue() -
a.getValue())} is not --- the subtraction can overflow and the cast then
reverses the order.}

% ------------------------------------------------------------------
\problem{Validate records and report the errors}{Medium}

Given records \texttt{"id,email,age"}, return one message per invalid record in
input order: \texttt{"line 3: bad age"}. A record is invalid if it has the
wrong field count, an email with no \texttt{@}, or an age outside 0 to 150.
Report only the first problem found in each record.

\idea{A validation ladder: check in a fixed order and return the first failure.
Keep the line number from the loop index, counting from 1.}

\solheading
\begin{lstlisting}[style=java]
public static List<String> validate(String[] rows) {
    List<String> errors = new ArrayList<>();
    for (int i = 0; i < rows.length; i++) {
        String problem = firstProblem(rows[i]);
        if (problem != null) errors.add("line " + (i + 1) + ": " + problem);
    }
    return errors;
}

private static String firstProblem(String row) {
    if (row == null) return "missing record";
    String[] parts = row.split(",", -1);
    if (parts.length != 3) return "wrong field count";
    if (!parts[1].contains("@")) return "bad email";
    int age;
    try { age = Integer.parseInt(parts[2].trim()); }
    catch (NumberFormatException e) { return "bad age"; }
    if (age < 0 || age > 150) return "bad age";
    return null;
}
\end{lstlisting}

\cost{$O(R)$.}
\trap{Line numbers are one-based in every message format anyone asks for, and
the loop index is zero-based. Off by one here fails every test at once, which
at least makes it easy to spot.}

% ------------------------------------------------------------------
\problem{Join two datasets on a key}{Medium}

One array holds \texttt{"userId,name"}, another holds
\texttt{"userId,orderTotal"}. Return \texttt{"name:total"} for every user
present in both, summing a user's orders, sorted by total descending.

\idea{Index the smaller side into a map, then stream the larger side against
it. This is a hash join, and it is $O(n+m)$ where the nested-loop version is
$O(nm)$.}

\solheading
\begin{lstlisting}[style=java]
public static List<String> join(String[] users, String[] orders) {
    Map<String,String> nameById = new HashMap<>();
    for (String u : users) {
        String[] p = u.split(",", -1);
        if (p.length == 2) nameById.put(p[0].trim(), p[1].trim());
    }

    Map<String,Long> totalById = new HashMap<>();
    for (String o : orders) {
        String[] p = o.split(",", -1);
        if (p.length != 2) continue;
        String id = p[0].trim();
        if (!nameById.containsKey(id)) continue;      // inner join
        try { totalById.merge(id, Long.parseLong(p[1].trim()), Long::sum); }
        catch (NumberFormatException ignored) { }
    }

    List<Map.Entry<String,Long>> rows = new ArrayList<>(totalById.entrySet());
    rows.sort((a, b) -> b.getValue().compareTo(a.getValue()));

    List<String> out = new ArrayList<>();
    for (Map.Entry<String,Long> e : rows)
        out.add(nameById.get(e.getKey()) + ":" + e.getValue());
    return out;
}
\end{lstlisting}

\cost{$O(n+m)$ plus the sort.}
\trap{Deciding what kind of join you were asked for. This one drops users
with no orders and orders with no user. A left join would keep every user
with a total of zero --- a different answer, and the statement will say
which.}

% ------------------------------------------------------------------
\problem{Bucket events by hour and find the peak}{Medium}

Each record is \texttt{"eventId,timestampSeconds"}. Count the events falling in
each hour and return the busiest hour as an hour index, breaking ties by the
earlier hour.

\idea{The bucket key is \texttt{timestamp / 3600} --- integer division is the
bucketing. Then a single scan for the maximum, taking \texttt{>} and not
\texttt{>=} so the earliest hour wins a tie.}

\solheading
\begin{lstlisting}[style=java]
public static long busiestHour(String[] events) {
    Map<Long,Integer> perHour = new HashMap<>();
    for (String e : events) {
        String[] p = e.split(",", -1);
        if (p.length != 2) continue;
        try {
            long hour = Long.parseLong(p[1].trim()) / 3600;
            perHour.merge(hour, 1, Integer::sum);
        } catch (NumberFormatException ignored) { }
    }
    if (perHour.isEmpty()) return -1;

    long bestHour = Long.MAX_VALUE;
    int bestCount = -1;
    for (Map.Entry<Long,Integer> e : perHour.entrySet()) {
        if (e.getValue() > bestCount
                || (e.getValue() == bestCount && e.getKey() < bestHour)) {
            bestCount = e.getValue();
            bestHour = e.getKey();
        }
    }
    return bestHour;
}
\end{lstlisting}

\cost{$O(R)$.}
\trap{Relying on iteration order to settle ties. \texttt{HashMap} order is
unspecified and changes with the keys, so the tie-break has to be written
out explicitly, as in the second clause above.}

% ------------------------------------------------------------------
\problem{Running balance per account}{Medium}

Transactions arrive as \texttt{"account,amount"} in chronological order.
Return, for each account, the lowest balance it ever reached.

\idea{One pass with two maps, or one map of a small array. Update the running
balance, then compare it to the stored minimum --- in that order, because the
minimum has to include the transaction you just applied.}

\solheading
\begin{lstlisting}[style=java]
public static Map<String,Long> lowestBalances(String[] transactions) {
    Map<String,Long> balance = new HashMap<>();
    Map<String,Long> lowest  = new HashMap<>();

    for (String t : transactions) {
        String[] p = t.split(",", -1);
        if (p.length != 2) continue;
        String account = p[0].trim();
        long amount;
        try { amount = Long.parseLong(p[1].trim()); }
        catch (NumberFormatException e) { continue; }

        long now = balance.getOrDefault(account, 0L) + amount;
        balance.put(account, now);
        lowest.merge(account, now, Math::min);
    }
    return lowest;
}
\end{lstlisting}

\cost{$O(R)$ time, $O(A)$ space for $A$ accounts.}
\trap{Seeding the minimum at \texttt{0} for every account. If an account only
ever receives money its lowest balance is its first deposit, not zero --- and
\texttt{merge} with no prior value stores the first balance, which is exactly
right.}
}{}
\IfFileExists{ch9-stretch.tex}{\chapterpage{Stretch}{Five problems above the bar. Nothing here is likely in a
seventy-minute screen for a final-year internship, and two of them were cut
from the assessment for exactly that reason. Work them when the first eight
chapters are comfortable --- the point is that a hard problem on the day
should feel familiar rather than alarming.}

% ------------------------------------------------------------------
\problem{Minimum travel time with k free roads}{Hard}

Cities $0 \ldots n-1$ are joined by bidirectional roads $\{u, v, t\}$. You hold
$k$ passes; using one on a road makes that road cost 0. Return the least time
from city 0 to city $n-1$, or $-1$.

\idea{The state you search is not a city, it is a \emph{pair}: the city, and how
many passes you have spent. Build $k+1$ copies of the graph --- layers --- and
run one Dijkstra over the lot. From $(u,p)$ each road either costs its time and
keeps you in layer $p$, or costs nothing and drops you into layer $p+1$.}

\solheading
\begin{lstlisting}[style=java]
public static int minTravelTime(int n, int[][] roads, int k) {
    if (n == 1) return 0;
    int layers = (int) Math.max(0L, Math.min((long) k, (long) n - 1));

    List<int[]>[] adj = new List[n];
    for (int i = 0; i < n; i++) adj[i] = new ArrayList<>();
    for (int[] r : roads) {
        adj[r[0]].add(new int[] {r[1], r[2]});
        if (r[0] != r[1]) adj[r[1]].add(new int[] {r[0], r[2]});
    }

    long[][] dist = new long[n][layers + 1];
    for (long[] row : dist) Arrays.fill(row, Long.MAX_VALUE);
    dist[0][0] = 0;

    // {cost, city, passesUsed}
    PriorityQueue<long[]> pq =
        new PriorityQueue<>((a, b) -> Long.compare(a[0], b[0]));
    pq.add(new long[] {0, 0, 0});

    while (!pq.isEmpty()) {
        long[] top = pq.poll();
        long cost = top[0];
        int u = (int) top[1], used = (int) top[2];
        if (cost > dist[u][used]) continue;          // stale entry
        if (u == n - 1) return (int) cost;

        for (int[] edge : adj[u]) {
            int v = edge[0];
            long paid = cost + edge[1];
            if (paid < dist[v][used]) {
                dist[v][used] = paid;
                pq.add(new long[] {paid, v, used});
            }
            if (used < layers && cost < dist[v][used + 1]) {
                dist[v][used + 1] = cost;
                pq.add(new long[] {cost, v, used + 1});
            }
        }
    }
    return -1;
}
\end{lstlisting}

\cost{$O(E \cdot K \log(V \cdot K))$ with $K = \min(k, n-1) + 1$.}
\trap{Three. $k$ can be $10^9$, so cap the layers at $n-1$ --- a simple path
uses no more roads than that, and without the cap you allocate a billion rows.
Distances overflow \texttt{int}, so accumulate in \texttt{long}. And the
tempting greedy --- find the shortest path, then zero its $k$ heaviest edges
--- is simply wrong: the best route when you have passes is often a different
route entirely.}
\alsoknown{``Cheapest flights with k stops'' and ``minimum cost path with k
refuels'' are the same layered-state idea.}

% ------------------------------------------------------------------
\problem{Expiring files}{Hard}

Extend the Part B \texttt{CloudStorage} so files can be given a lifetime.
\texttt{addFileAt(timestamp, user, name, size, ttl)} creates a file alive over
$[\textit{timestamp}, \textit{timestamp}+\textit{ttl})$;
\texttt{getFileSizeAt}, \texttt{deleteFileAt} and \texttt{getNLargestAt} see
only files alive at their timestamp. An expired file frees its owner's
capacity.

\idea{The level attacks the incremental \texttt{used} counter that made the
earlier levels cheap: bytes now free themselves with no operation to hook.
Rather than recompute usage on every query, make expiry \emph{lazy} --- each
timestamped call first deletes everything already dead at its timestamp,
through the ordinary delete path, and then delegates to the untimed method
that already exists. One notion of ``the files that exist'', and every earlier
level untouched.}

\solheading
\begin{lstlisting}[style=java]
// FileEntry gains: final long expiresAt;   (Long.MAX_VALUE = never)

public Integer addFileAt(int timestamp, String user,
                         String name, int size, int ttl) {
    if (ttl <= 0 || timestamp < 0) return null;
    purge(timestamp);
    // long, because timestamp + ttl overflows int near Integer.MAX_VALUE
    return addEntry(user, name, size, (long) timestamp + (long) ttl);
}

public Integer getFileSizeAt(int timestamp, String name) {
    purge(timestamp);
    return getFileSize(name);
}

/** Drops every file dead at `now`, crediting each owner's capacity. */
private void purge(long now) {
    List<String> expired = null;
    for (Map.Entry<String,FileEntry> e : files.entrySet())
        if (e.getValue().expiresAt <= now) {
            if (expired == null) expired = new ArrayList<>();
            expired.add(e.getKey());
        }
    if (expired == null) return;
    for (String name : expired) removeEntry(name);   // credits the owner
}
\end{lstlisting}

\cost{$O(F)$ per timestamped call in the worst case, $O(1)$ amortised per file
since each file is purged once.}
\trap{Collecting the dead keys before deleting --- removing from
\texttt{files} while iterating it throws
\texttt{ConcurrentModificationException}. Storing an \emph{absolute} expiry
rather than a TTL, so backup and restore do not restart a lifetime. And
treating the interval as closed: a file with \texttt{ttl} 10 created at 0 is
dead \emph{at} 10, not after it.}

% ------------------------------------------------------------------
\problem{LRU cache}{Hard}

A fixed-capacity cache with $O(1)$ \texttt{get} and \texttt{put}. When it is
full, the least recently used entry is evicted. A \texttt{get} counts as a use.

\idea{\texttt{LinkedHashMap} does this in a dozen lines: construct it in
access-order mode and override \texttt{removeEldestEntry}. Know the manual
version too --- a hash map to nodes plus a doubly linked list --- because
interviewers sometimes ban the shortcut.}

\solheading
\begin{lstlisting}[style=java]
class LruCache<K, V> extends LinkedHashMap<K, V> {
    private final int capacity;

    LruCache(int capacity) {
        // initialCapacity, loadFactor, accessOrder=true
        super(16, 0.75f, true);
        this.capacity = capacity;
    }

    @Override
    protected boolean removeEldestEntry(Map.Entry<K, V> eldest) {
        return size() > capacity;
    }
}
\end{lstlisting}

\cost{$O(1)$ expected for both operations.}
\trap{The third constructor argument. Left at \texttt{false} --- the default
--- you get insertion order, which evicts the oldest \emph{inserted} entry and
makes it an FIFO cache, not an LRU one. The bug passes any test that never
re-reads a key.}

% ------------------------------------------------------------------
\problem{Coin change}{Hard}

Given coin denominations and a target, return the fewest coins that sum to it,
or $-1$. Each denomination may be used any number of times.

\idea{Bottom-up dynamic programming over the amount. The best way to make
$a$ is one coin on top of the best way to make $a - \textit{coin}$, minimised
over the coins. Fill the table upwards so every subproblem is already solved
when you need it.}

\solheading
\begin{lstlisting}[style=java]
public static int coinChange(int[] coins, int amount) {
    if (amount < 0) return -1;
    int[] best = new int[amount + 1];
    Arrays.fill(best, Integer.MAX_VALUE);
    best[0] = 0;                               // zero coins make zero

    for (int a = 1; a <= amount; a++)
        for (int coin : coins) {
            if (coin <= 0 || coin > a) continue;
            if (best[a - coin] == Integer.MAX_VALUE) continue;
            best[a] = Math.min(best[a], best[a - coin] + 1);
        }

    return best[amount] == Integer.MAX_VALUE ? -1 : best[amount];
}
\end{lstlisting}

\cost{$O(\textit{amount} \times \textit{coins})$ time, $O(\textit{amount})$
space.}
\trap{The greedy ``take the biggest coin that fits'' is wrong on coin sets
like $\{1, 3, 4\}$, where making 6 greedily takes $4+1+1$ and the answer is
$3+3$. It happens to work on real currency, which is why people trust it. Also
guard \texttt{best[a - coin] == MAX\_VALUE} before adding one, or it
overflows to a negative and poisons the table.}

% ------------------------------------------------------------------
\problem{Longest increasing subsequence}{Hard}

Return the length of the longest strictly increasing subsequence --- the
elements need not be adjacent. \texttt{[10,9,2,5,3,7,101,18]} gives 4.

\idea{Keep a list \texttt{tails}, where \texttt{tails[i]} is the smallest value
that can end an increasing subsequence of length $i+1$. That list is always
sorted, so each new element is placed by binary search: it either extends the
list or lowers an existing tail, making future extensions easier. The list is
not itself a valid subsequence --- only its \emph{length} is the answer.}

\solheading
\begin{lstlisting}[style=java]
public static int lengthOfLIS(int[] a) {
    if (a == null || a.length == 0) return 0;
    int[] tails = new int[a.length];
    int size = 0;

    for (int x : a) {
        int lo = 0, hi = size;
        while (lo < hi) {                      // first tail >= x
            int mid = (lo + hi) >>> 1;
            if (tails[mid] < x) lo = mid + 1;
            else hi = mid;
        }
        tails[lo] = x;
        if (lo == size) size++;
    }
    return size;
}
\end{lstlisting}

\cost{$O(n\log n)$ time, $O(n)$ space. The straightforward DP is $O(n^2)$ and
is the right answer below a few thousand elements.}
\trap{\texttt{(lo + hi) / 2} overflows for huge bounds; \texttt{>>> 1} is the
habit worth having. And for a \emph{non-strictly} increasing subsequence the
binary search changes to ``first tail $>$ x'' --- a one-character difference
that silently changes the answer.}
}{}

\end{document}
